\documentclass[profile=standard,twoside=false]{ipara-handbook}

\usetikzlibrary{arrows.meta,calc,fit,positioning,backgrounds}
\definecolor{themegreen}{HTML}{4ade80}
\pgfdeclarelayer{background}
\pgfdeclarelayer{foreground}
\pgfsetlayers{background,main,foreground}

\handbooksetup
  {AI-I04｜现代 LLM 应用系统：Context、Retrieval、RAG、Memory、Agent、Workflow 与 MCP}
  {人工智能 · Integration · AI-I04}
  {Context · Retrieval · Memory · Agent · MCP}

\begin{document}

\begin{titlepage}
  \centering
  \vspace*{0.16\textheight}
  {\Huge\bfseries\sffamily AI-I04｜现代 LLM 应用系统\par}
  \vspace{0.8em}
  {\Large Context、Retrieval、RAG、Memory、Agent、Workflow 与 MCP\par}
  \vspace{1.6em}
  {\small 技术快照：2026-08-22\par}
  \vspace{1.1em}
  {\small\href{../../40_复试/10_人工智能与机器学习/90_综合专题/AI-I04_现代LLM应用系统/README.md}{返回 Integration Landing Page}\par}
\end{titlepage}

\tableofcontents

\section{本册定位：从语言模型到可运行系统}

一个语言模型最基本的运行形式是“输入当前上下文，生成输出”。真实应用却还要处理私有知识、实时业务数据、工具调用、多步状态、权限、失败恢复和质量评估。

本册不重新解释 Transformer、概率、优化或搜索等单点机制，只追踪这些成熟模块怎样在一次真实任务里交接。母问题是：

\begin{mentalmodel}{Mother Question｜一次 LLM 任务到底怎样跑完}
一次真实 LLM 任务执行时，哪些信息进入当前上下文，哪些知识需要检索，什么时候只生成文本，什么时候调用工具，谁控制多步流程，状态保存在哪里，外部系统怎样安全接入，最后又怎样判断这次执行是否真的成功？
\end{mentalmodel}

最小运行链可以写成：

\begin{processchain}
  \processstage{用户任务}
  \processstage{Context 组装}
  \processstage{模型判断}
  \processstage{生成或调用工具}
  \processstage{状态更新}
  \processstage{验证与结束}
\end{processchain}

这条链只表示一次任务的运行顺序，不表示各阶段在理论上属于同一个机制。

\begin{handbookdiagram}
\centering
\begin{tikzpicture}[
  >=Stealth,
  font=\small,
  color=themefg,
  text=themefg,
  main/.style={draw=themecurve,fill=themecurve!14!themebg,rounded corners=4pt,minimum width=2.7cm,minimum height=0.9cm,align=center,line width=1pt},
  support/.style={draw=themeamber,fill=themeamber!13!themebg,rounded corners=4pt,minimum width=2.55cm,minimum height=0.86cm,align=center,line width=0.95pt},
  external/.style={draw=themegreen,fill=themegreen!12!themebg,rounded corners=4pt,minimum width=2.65cm,minimum height=0.86cm,align=center,line width=0.95pt},
  guard/.style={draw=themealert,fill=themealert!10!themebg,rounded corners=4pt,minimum width=2.75cm,minimum height=0.9cm,align=center,line width=0.95pt},
  flow/.style={->,draw=themefg!90,line width=1.05pt},
  retrieval/.style={->,draw=themeamber,line width=1pt,dashed},
  action/.style={->,draw=themegreen,line width=1pt},
  feedback/.style={->,draw=themecurve,line width=1pt,dashed}
]

% Three semantic regions. Horizontal position expresses responsibility, not time.
\begin{pgfonlayer}{background}
  \fill[fill=themeamber!5!themebg,draw=themeamber!45,dashed,rounded corners=7pt]
    (-0.2,1.45) rectangle (4.0,9.7);
  \fill[fill=themecurve!5!themebg,draw=themecurve!45,dashed,rounded corners=7pt]
    (4.35,0.25) rectangle (10.7,10.35);
  \fill[fill=themegreen!5!themebg,draw=themegreen!45,dashed,rounded corners=7pt]
    (11.05,1.45) rectangle (15.75,9.7);
\end{pgfonlayer}

\node[anchor=west,font=\bfseries,color=themeamber] at (0.05,9.25) {外部知识};
\node[anchor=west,font=\bfseries,color=themecurve] at (4.65,9.9) {应用运行时};
\node[anchor=west,font=\bfseries,color=themegreen] at (11.35,9.25) {外部能力};

% Knowledge path: one-way preprocessing and retrieval support.
\node[support] (docs)  at (1.9,7.8) {文档 / 文件};
\node[support] (chunk) at (1.9,6.15) {解析与 Chunk};
\node[support] (embed) at (1.9,4.5) {Embedding};
\node[support] (index) at (1.9,2.85) {Retrieval Index};
\draw[flow,draw=themeamber] (docs) -- (chunk);
\draw[flow,draw=themeamber] (chunk) -- (embed);
\draw[flow,draw=themeamber] (embed) -- (index);

% Runtime primary path.
\node[main] (user)    at (7.5,8.55) {用户任务};
\node[main] (context) at (7.5,6.85) {Context Assembly};
\node[main] (model)   at (7.5,5.05) {LLM / Model Decision};
\node[main] (answer)  at (5.75,2.75) {最终输出};
\node[guard] (runtime) at (9.25,2.75) {Runtime\\\scriptsize 校验 · 路由 · 授权};
\node[main] (state)   at (7.5,0.95) {State / Checkpoint};
\draw[flow] (user) -- (context);
\draw[flow] (context) -- (model);
\draw[flow] (model) -| node[pos=0.35,above=1pt,font=\scriptsize]{直接回答} (answer);
\draw[flow,draw=themealert] (model) -| node[pos=0.35,above=1pt,font=\scriptsize,color=themealert]{Tool Call} (runtime);

% Retrieval injects evidence only into context assembly.
\draw[retrieval] (index.east) -- ++(1.1,0) |- node[pos=0.7,above=1pt,font=\scriptsize,color=themeamber]{检索证据} (context.west);

% External capability path.
\node[external] (tool)    at (13.35,6.25) {Tool / MCP Client};
\node[external] (server)  at (13.35,4.45) {MCP Server / API};
\node[external] (backend) at (13.35,2.65) {业务系统\\\scriptsize DB · SaaS · Service};
\draw[action] (runtime.east) -- ++(0.8,0) |- node[pos=0.68,right=2pt,font=\scriptsize,color=themegreen]{执行} (tool.south);
\draw[action] (tool) -- (server);
\draw[action] (server) -- (backend);

% Observation updates state; the next turn re-enters context assembly via an outer corridor.
\draw[feedback] (backend.west) -- ++(-1.2,0) |- node[pos=0.72,below=1pt,font=\scriptsize,color=themecurve]{Observation} (state.east);
\draw[feedback] (state.east) -- ++(6.6,0) |- node[pos=0.6,right=2pt,font=\scriptsize,color=themecurve]{下一轮} (context.east);

\node[anchor=west,font=\scriptsize\bfseries,color=themealert] at (4.65,0.25)
  {贯穿全程：Security · Evaluation · Observability · Cost / Latency · Human Approval · Recovery};

\end{tikzpicture}

\diagramcaption{现代 LLM 应用系统总架构。主路径表示一次任务怎样推进；左侧是外部知识进入 Context 的检索路径，右侧是模型通过 Tool / MCP 访问外部能力的执行路径。}
\end{handbookdiagram}

\subsection{本册与各 Core Area 的 Ownership}

\begin{handbooktable}{L{0.27\linewidth}L{0.28\linewidth}Y}
\toprule
\textbf{对象 / 机制} & \textbf{Canonical Owner} & \textbf{AI-I04 只负责什么}\\
\midrule
Transformer / Attention & Area 60 深度学习 & 模型在完整系统中承担生成与推理运行时\\
Autoregressive generation & Area 80 生成模型 & 生成结果怎样接入应用流程\\
Learning / evaluation contract & Area 40 机器学习 & 系统级评测怎样拆成 retrieval、answer、trajectory\\
Optimization / AutoDiff & Area 50 优化与学习计算 & 本册只在成本或调用边界处引用\\
Search / decoding search & Area 10 问题求解 & Agent 或生成流程何时调用搜索机制\\
Decision / sequential control & Area 70 决策、控制与 RL & 多步行动和策略语义的接口\\
Context / Retrieval / RAG / Tool / Workflow / MCP 的组合 & 本 Integration & 一次任务中对象怎样传递、状态怎样变化、控制权在哪里\\
\bottomrule
\end{handbooktable}

必须稳定区分：

\begin{boundarytable}
\boundarytablehead
Vector Store & $\neq$ & Long-term Memory & Vector Store 是检索索引；Memory 是跨轮次或跨任务持久保存的状态/事实 & 把“能检索”误当成“模型自动记住”\\
Retrieval & $\neq$ & RAG & Retrieval 只负责找证据；RAG 还要把证据组织进 Context 并生成回答 & 只评向量搜索却声称整个 RAG 有效\\
Tool Calling & $\neq$ & Agent & Tool Calling 是一次结构化动作接口；Agent 还需要循环、状态、停止条件和控制策略 & 把会调函数的单次调用误称为完整 Agent\\
MCP & $\neq$ & Agent & MCP 标准化外部 Context / Capability 接口；Agent 决定任务怎样推进 & 把协议本身当成“智能”或“安全机制”\\
Graph State & $\neq$ & Long-term Memory & State 服务当前 workflow；Long-term Memory 跨 thread / session 持久存在 & 状态越堆越大，生命周期和权限不清\\
\bottomrule
\end{boundarytable}

\section{Context：模型这一轮真正看见什么}

上下文窗口（context window）表示一次模型调用能够处理的 token 工作集。它可以包含系统/开发者指令、用户输入、对话历史、检索到的文档、工具定义、工具结果、文件和任务状态。

但下面这些对象不是同一个东西：

\begin{itemize}
  \item \textbf{预训练知识}：模型参数中已经学到的统计规律；
  \item \textbf{当前 Context}：这一轮请求实际送给模型的工作集；
  \item \textbf{Retrieval Index}：外部知识的可检索索引；
  \item \textbf{Application Database}：业务系统的事实与事务状态；
  \item \textbf{Long-term Memory}：应用显式保存、以后任务可以再次读取的信息。
\end{itemize}

截至 2026-08-22，OpenAI GPT-5.6 Sol / Terra / Luna 的 context window 都是约 $1.05$M tokens，最大输出为 $128$K tokens。长上下文扩大了“能放多少”的上限，但没有自动解决“该放什么”。

\begin{mechanism}{Context Engineering｜真正优化的是相关而充分的工作集}
Context 越长并不天然越好。无关历史、重复资料和不可信外部文本都会增加干扰、成本和延迟。系统要做的是从“所有可用信息”中选择“当前任务需要的信息”，再把它按清楚的角色和边界组织给模型。
\end{mechanism}

每轮调用至少要回答：

\begin{enumerate}
  \item 哪些 instructions 必须保留，优先级怎样；
  \item 哪些历史消息仍然相关；
  \item 哪些知识应该实时 Retrieval，而不是长期堆在对话里；
  \item 哪些 Tool schemas 需要暴露给模型；
  \item Tool result 哪部分要进入下一轮；
  \item 哪些状态应该结构化保存，而不是继续扩充消息历史；
  \item 哪些旧信息已经过期；
  \item 哪些外部内容必须按不可信数据处理。
\end{enumerate}

因此 Prompt Engineering 是 Context Engineering 的一部分，但不是全部。

\section{Retrieval：从外部知识中选择当前证据}

\subsection{Embedding 是表示，不是知识库本身}

文本 Embedding 可以写成：

$$
f_{\text{embed}}(x)=\boldsymbol{v}\in\mathbb{R}^{d}.
$$

维度 $d$ 由具体 embedding model 决定，不存在统一固定值。Embedding 的作用是把文本或其他对象映射到适合相似度计算、聚类和检索的 learned representation。

典型语义检索过程是：

\begin{processchain}
  \processstage{Document Chunk}
  \processstage{Embedding}
  \processstage{Vector Index}
  \processstage{Query Embedding}
  \processstage{Similarity / ANN}
  \processstage{Candidate Chunks}
\end{processchain}

实际相似度可以使用 cosine、dot product、Euclidean distance 等，具体取决于 embedding model 和索引设计。

截至当前版本，OpenAI 仍提供 \codepath{text-embedding-3-small} 与 \codepath{text-embedding-3-large}。\codepath{text-embedding-3} 系列支持通过 \code{dimensions} 调整输出维度。轻量本地教学场景中，也常见 \codepath{sentence-transformers/all-MiniLM-L6-v2}，它输出 $384$ 维向量。具体型号只是实现选择，不能反过来定义 Embedding 的机制。

\subsection{Chunking：检索单元必须服从文档结构}

长文档通常不适合作为一个 retrieval unit。Chunk 太大时，一个向量可能混入多个主题；Chunk 太小时，定义、条件或论证链可能被切断。

Overlap 可以保护边界附近的上下文，但不存在跨数据集通用的“最佳 chunk size / overlap”。更稳的判断顺序是：

\begin{enumerate}
  \item 先看文档原生结构；
  \item 再决定一个可独立理解的语义单元；
  \item 再看真实 query 通常需要多大证据范围；
  \item 最后结合 Context budget 与 retrieval eval 调参数。
\end{enumerate}

例如：

\begin{itemize}
  \item 法律文本优先保留条、款、定义和引用关系；
  \item API 文档尽量保持函数、参数、返回值完整；
  \item 对话记录可以细分，但要保留说话者和必要邻近轮次；
  \item 教材不应机械切断定义、定理、证明和例题的局部闭环。
\end{itemize}

OpenAI 当前 Vector Store 的 \code{auto} chunking 默认值为：

\begin{lstlisting}
max_chunk_size_tokens = 800
chunk_overlap_tokens  = 400
\end{lstlisting}

这是当前产品默认值，不是理论最优值。生产系统仍然要用自己的查询集做 retrieval eval。

\subsection{Retrieval 是一条筛选链，不只是最近邻搜索}

完整 retrieval pipeline 通常包含：

\begin{processchain}
  \processstage{Query Rewrite}
  \processstage{Lexical / Vector Search}
  \processstage{Metadata / Permission Filter}
  \processstage{Candidate Retrieval}
  \processstage{Rerank}
  \processstage{Top-k / Threshold}
\end{processchain}

当前 OpenAI Vector Store Search 已显式提供 attribute filters、\codepath{max_num_results}、ranker、\codepath{score_threshold} 和 query rewrite 等能力。这说明检索真正要解决的是“候选怎么生成、怎样过滤、怎样排序、最终哪些证据进入 Context”。

若接受阈值为 $\tau$，提高 $\tau$ 往往会提高 precision，同时可能降低 recall。阈值必须由任务风险和测试集决定，不能只凭感觉指定。

SQL 与 Vector Retrieval 也不是二选一。同一个系统可以同时使用结构化过滤、全文搜索、向量相似度、权限过滤与 reranker；例如 PostgreSQL 配合 pgvector 就能在一个数据系统中完成结构化查询和向量近邻搜索。

\begin{criterion}{先判断问题需要哪一种“真相接口”}
不要因为系统里有 Vector Store，就把所有数据都改写成 Embedding。先看用户问题要求什么类型的答案：

\begin{handbooktable}{L{0.24\linewidth}L{0.28\linewidth}Y}
\toprule
\textbf{问题对象} & \textbf{优先接口} & \textbf{为什么}\\
\midrule
精确结构化事实 & SQL / typed API / RPC & ID、价格、库存、状态、时间范围、聚合和权限过滤需要确定语义与精确条件\\
长文本中的语义证据 & FTS / Vector / Hybrid Retrieval & 用户表述与原文未必共享关键词，需要按语义召回相关段落\\
实时外部事实 & Tool / API & 当前值属于外部系统事实，应在请求时读取，而不是依赖旧 Context 或旧 embedding\\
有副作用的动作 & Authorized Tool / Workflow & 发邮件、退款、写数据库等必须经过 schema、权限、幂等和审计边界\\
\bottomrule
\end{handbooktable}
\end{criterion}

一个 Agentic RAG 系统真正“agentic”的地方，不是 Vector Store 本身，而是 Runtime 根据当前意图和证据需求，在这些接口之间做路由。例如“融资政策是什么”可能需要从政策文档做语义检索；“当前有哪些蓝色且低于 3 万美元的库存”应优先查询结构化库存表；“把其中一辆加入试驾日程”才进入有副作用的 Tool / Workflow。三类问题可以出现在同一次对话中，但不能共用同一种数据访问语义。

\section{RAG：把外部证据送入当前生成过程}

RAG（Retrieval-Augmented Generation）的责任链是：

\begin{enumerate}
  \item \textbf{Retrieval}：当前问题需要哪些外部证据；
  \item \textbf{Augmentation}：怎样把证据、来源和用户问题组织进当前 Context；
  \item \textbf{Generation}：模型怎样基于这组证据生成回答。
\end{enumerate}

因此：

$$
\boxed{\text{RAG}\neq\text{Vector Search}}
$$

Vector Search 只是 Retrieval 的一种实现。RAG 还拥有 request-time context construction 与 generation 这两段组合责任。

\begin{handbookdiagram}
\centering
\begin{tikzpicture}[
  >=Stealth,
  font=\small,
  color=themefg,
  text=themefg,
  idx/.style={draw=themeamber,fill=themeamber!13!themebg,rounded corners=4pt,minimum width=2.25cm,minimum height=0.86cm,align=center,line width=0.95pt},
  run/.style={draw=themecurve,fill=themecurve!13!themebg,rounded corners=4pt,minimum width=2.25cm,minimum height=0.86cm,align=center,line width=0.95pt},
  evidence/.style={draw=themegreen,fill=themegreen!12!themebg,rounded corners=4pt,minimum width=2.35cm,minimum height=0.86cm,align=center,line width=0.95pt},
  indexflow/.style={->,draw=themeamber,line width=1pt},
  requestflow/.style={->,draw=themecurve,line width=1pt},
  evidencejoin/.style={->,draw=themegreen,line width=1pt,dashed}
]

\begin{pgfonlayer}{background}
  \fill[fill=themeamber!5!themebg,draw=themeamber!45,dashed,rounded corners=7pt]
    (-0.15,4.35) rectangle (14.2,7.35);
  \fill[fill=themecurve!5!themebg,draw=themecurve!45,dashed,rounded corners=7pt]
    (-0.15,0.15) rectangle (16.0,4.0);
\end{pgfonlayer}

\node[anchor=west,font=\bfseries,color=themeamber] at (0.05,6.95) {索引路径：把外部知识整理成可检索结构};
\node[anchor=west,font=\bfseries,color=themecurve] at (0.05,3.62) {请求路径：为当前问题选择证据，再把证据送入模型};

% Index path.
\node[idx] (source) at (1.35,5.55) {Document Source};
\node[idx] (chunk)  at (3.9,5.55) {Parse / Chunk};
\node[idx] (embed)  at (6.45,5.55) {Embedding};
\node[idx] (index)  at (9.0,5.55) {Retrieval Index};
\draw[indexflow] (source) -- (chunk);
\draw[indexflow] (chunk) -- (embed);
\draw[indexflow] (embed) -- (index);
\node[anchor=west,font=\scriptsize,color=themegray] at (10.35,5.55) {content reference + metadata + permissions};

% Request path. Each edge is request-time data/control flow.
\node[run]      (query)    at (1.35,2.35) {User Query};
\node[run]      (retrieve) at (3.9,2.35) {Retrieve\\\scriptsize rewrite / filter};
\node[evidence] (select)   at (6.45,2.35) {Evidence Selection\\\scriptsize rerank / top-k / threshold};
\node[evidence] (context)  at (9.15,2.35) {Context Assembly};
\node[run]      (model)    at (11.85,2.35) {LLM};
\node[run]      (answer)   at (14.55,2.35) {Answer / Citation};
\draw[requestflow] (query) -- (retrieve);
\draw[requestflow] (retrieve) -- (select);
\draw[requestflow] (select) -- (context);
\draw[requestflow] (context) -- (model);
\draw[requestflow] (model) -- (answer);

% The two paths meet only here: the index serves retrieval.
\draw[evidencejoin] (index.south) -- ++(0,-0.75) -| node[pos=0.72,above=1pt,font=\scriptsize,color=themegreen]{候选内容} (retrieve.north);

\node[anchor=west,font=\scriptsize,color=themegray] at (0.1,0.55)
  {Retrieval 找证据；RAG 还包括 Context Assembly 与 Generation。两条路径不能压成“向量库 = RAG”。};

\end{tikzpicture}

\diagramcaption{RAG 的两条路径。上半部分是离线/增量索引路径，下半部分是请求时路径；二者在 Retrieval Index 处汇合，检索出的证据经过 Context Assembly 才进入模型。}
\end{handbookdiagram}

\subsection{RAG 与 Fine-tuning 的边界}

Fine-tuning 改变模型参数或行为分布；RAG 不修改基础模型权重，而是在 inference time 提供外部证据。频繁更新的企业知识、需要 citation 的问答、私有文档通常优先考虑 Retrieval / RAG；行为风格、稳定格式和高频任务模式则可能更适合 fine-tuning。两者也可以组合。

\subsection{Grounded Generation 仍然可能失败}

即使系统要求“只根据提供的文档回答”，仍然可能出现：

\begin{itemize}
  \item Retrieval miss：正确证据没有进入 Context；
  \item Retrieval false positive：检索了看似相似但实际不相关的内容；
  \item Stale source：证据已经过期；
  \item Conflicting sources：多个来源互相冲突；
  \item Incorrect synthesis：证据正确，但模型组合或解释错误；
  \item Citation mismatch：回答与引用证据并不真正对应。
\end{itemize}

因此 RAG Evaluation 至少要拆成 Retrieval Quality 与 Generation / Grounding Quality 两层，不能只用一个“回答准确率”覆盖全部失败原因。

\section{Model API：当前应用层怎样调用模型}

截至 2026-08-22，OpenAI 当前文档建议使用 Responses API 构建 reasoning、tool-calling 与 multi-turn workflow。最小 Python 调用可以写成：

\begin{lstlisting}[language=Python]
from openai import OpenAI

client = OpenAI()
response = client.responses.create(
    model="gpt-5.6",
    input="Explain retrieval-augmented generation in one paragraph.",
)
print(response.output_text)
\end{lstlisting}

真正应该记的是调用契约，而不是某一版 response object 的路径：

\begin{processchain}
  \processstage{Client}
  \processstage{Authentication}
  \processstage{Model + Input}
  \processstage{Instructions + Tools}
  \processstage{Response}
  \processstage{Output / Usage / Events}
\end{processchain}

Responses API 当前可以组合 function tools、web search、file search、code interpreter、computer use、remote MCP 等能力。模型调用已经不再只是“一问一答”，而可以成为受控工具运行时的一部分。

\subsection{模型选择是一项系统决策}

GPT-5.6 family 当前包括：

\begin{itemize}
  \item \code{gpt-5.6-sol} / \code{gpt-5.6}：能力优先；
  \item \code{gpt-5.6-terra}：能力与成本平衡；
  \item \code{gpt-5.6-luna}：高吞吐、成本敏感。
\end{itemize}

三者当前 context window 均约 $1.05$M tokens。选择模型时不能只看公开 benchmark，还要在自己的任务上比较 Task Success、Quality、Latency、Token Usage、Tool-call Accuracy 与 Cost。

\section{Tool Calling：模型提出动作，Runtime 负责执行}

Tool 至少应明确：

\begin{itemize}
  \item name 与 description；
  \item input schema；
  \item output contract；
  \item permission boundary；
  \item side-effect semantics；
  \item failure contract。
\end{itemize}

一次 Tool Call 的责任链是：

\begin{processchain}
  \processstage{Model 提出 Tool Call}
  \processstage{Runtime 校验参数}
  \processstage{鉴权 / 用户确认}
  \processstage{执行 Tool}
  \processstage{返回 Observation}
  \processstage{进入下一轮}
\end{processchain}

所以“模型调用工具”并不意味着模型自己执行了外部副作用。真正执行数据库写入、发邮件、调用业务 API 的仍然是应用 Runtime 或 Tool Server。

\section{Agent：核心是受控循环}

Agent 不是“会用工具的更聪明 LLM”。它至少包含模型决策、工具系统、状态、循环与停止条件。

一次任务可以还原为：

\begin{enumerate}
  \item 接收用户任务；
  \item 组装当前 Context；
  \item 模型选择回答、Tool Call、route / handoff 或 stop；
  \item Runtime 校验动作；
  \item Tool 在权限、timeout、retry 和副作用规则下执行；
  \item Observation 返回；
  \item 更新 conversation / task state / checkpoint / trace；
  \item 继续下一轮或终止；
  \item 高风险动作必要时插入人工批准；
  \item 全程记录 token、latency、cost 与 error。
\end{enumerate}

\begin{handbookdiagram}
\centering
\begin{tikzpicture}[
  >=Stealth,
  font=\small,
  color=themefg,
  text=themefg,
  core/.style={draw=themecurve,fill=themecurve!13!themebg,rounded corners=4pt,minimum width=2.8cm,minimum height=0.9cm,align=center,line width=1pt},
  guard/.style={draw=themealert,fill=themealert!10!themebg,rounded corners=4pt,minimum width=2.8cm,minimum height=0.9cm,align=center,line width=1pt},
  tool/.style={draw=themegreen,fill=themegreen!12!themebg,rounded corners=4pt,minimum width=2.8cm,minimum height=0.9cm,align=center,line width=1pt},
  state/.style={draw=themeamber,fill=themeamber!12!themebg,rounded corners=4pt,minimum width=2.8cm,minimum height=0.9cm,align=center,line width=1pt},
  flow/.style={->,draw=themefg!90,line width=1.05pt},
  conditional/.style={->,draw=themealert,line width=1pt,dashed},
  feedback/.style={->,draw=themeamber,line width=1pt,dashed},
  support/.style={->,draw=themegray,line width=0.95pt,densely dashed}
]

% Primary loop is vertical. Side branches never cross it.
\node[core]  (context)  at (6.5,8.1) {Context Assembly};
\node[core]  (model)    at (6.5,6.25) {Model Decision};
\node[guard] (validate) at (6.5,4.35) {Validate / Authorize};
\node[tool]  (execute)  at (6.5,2.45) {Tool Execution};
\node[state] (observe)  at (6.5,0.55) {Observation + State Update};
\draw[flow] (context) -- (model);
\draw[flow,draw=themealert] (model) -- node[right=2pt,font=\scriptsize,color=themealert]{Tool Call} (validate);
\draw[flow,draw=themegreen] (validate) -- node[right=2pt,font=\scriptsize,color=themegreen]{允许执行} (execute);
\draw[flow,draw=themegreen] (execute) -- (observe);

% Termination branch: answer is a terminal output, not another loop state.
\node[core,minimum width=2.65cm] (answer) at (2.25,6.25) {Final Answer\\\scriptsize terminate};
\draw[flow] (model.west) -- node[above=1pt,font=\scriptsize]{直接回答} (answer.east);

% Human approval is conditional and only gates sensitive tool calls.
\node[guard,minimum width=2.75cm] (human) at (10.75,3.4) {Human Approval\\\scriptsize high-risk only};
\draw[conditional] (validate.east) -| node[pos=0.45,above=1pt,font=\scriptsize,color=themealert]{敏感操作} (human.north);
\draw[conditional] (human.south) |- node[pos=0.55,below=1pt,font=\scriptsize,color=themealert]{批准 / 拒绝} (execute.east);

% Checkpoint is a persistence branch. It does not decide the next action.
\node[state,minimum width=2.75cm] (checkpoint) at (10.75,0.55) {Checkpoint / Trace};
\draw[support] (observe.east) -- (checkpoint.west);

% Normal next turn returns through the left outer corridor.
\draw[feedback] (observe.west) -- ++(-4.7,0) |- node[pos=0.62,left=2pt,font=\scriptsize,color=themeamber]{下一轮} (context.west);

% Recovery / replay returns from persisted state through the far-right corridor.
\draw[support] (checkpoint.east) -- ++(1.55,0) |- node[pos=0.58,right=2pt,font=\scriptsize,color=themegray]{恢复 / 回放} (context.east);

\node[anchor=west,font=\scriptsize,color=themegray] at (0.15,9.05)
  {主循环只回答“下一步怎样执行”；人工批准和 checkpoint 都是旁路控制，不能画成模型自主性的组成部分。};

\end{tikzpicture}

\diagramcaption{Agent 受控运行时循环。模型负责提出下一步候选动作，Runtime 负责验证与执行；高风险操作可以在执行前进入人工批准，Observation 和 State 再回到下一轮。}
\end{handbookdiagram}

\subsection{Planning 与 Stopping 必须显式设计}

一个真实 Agent 至少要回答：

\begin{itemize}
  \item 谁决定下一步；
  \item 最多运行多少轮；
  \item 什么条件算任务完成；
  \item 什么时候因为证据不足停止；
  \item 什么时候请求用户确认；
  \item 什么时候转给人工；
  \item 怎样避免循环和重复副作用。
\end{itemize}

Agent 也不必什么都负责。Routing、Planning、Retrieval、Tool Selection、Execution、Verification、Summarization 与 Escalation 可以分别由模型、确定性代码、workflow edge 或人工拥有。

\section{Agent Harness：把一个模型变成可长期运行的代理系统}

“Agent Harness”不是模型能力的同义词，也没有一个跨所有厂商完全统一的字段清单。这里采用当前工程实践中最有操作性的定义：\textbf{Harness 是围绕模型建立的一组运行时约束与支撑机制，使模型能够在明确的上下文、工具、状态、权限、循环和观测边界内持续完成任务。}

可以把它压成：

$$
\boxed{
\text{Agent Harness}
=
\text{Context}
+\text{Tools}
+\text{Loop / Planning}
+\text{Session / State}
+\text{Memory / Artifacts}
+\text{Policy / Permissions}
+\text{Trace / Eval Hooks}
}
$$

模型仍然负责根据当前输入产生候选文本或动作；Harness 决定模型这一轮看到什么、能调用什么、动作能否执行、怎样继续下一轮、状态怎样持久化以及失败后怎样恢复。换一个更强的模型，并不会自动替你完成这些系统责任。

\begin{handbookdiagram}
\centering
\begin{tikzpicture}[
  >=Stealth,
  font=\small,
  color=themefg,
  text=themefg,
  model/.style={draw=themecurve,fill=themecurve!16!themebg,rounded corners=5pt,minimum width=3.05cm,minimum height=1.0cm,align=center,line width=1.05pt},
  runtime/.style={draw=themeamber,fill=themeamber!11!themebg,rounded corners=4pt,minimum width=2.8cm,minimum height=0.86cm,align=center,line width=0.95pt},
  guard/.style={draw=themealert,fill=themealert!10!themebg,rounded corners=4pt,minimum width=2.9cm,minimum height=0.86cm,align=center,line width=0.95pt},
  state/.style={draw=themegreen,fill=themegreen!10!themebg,rounded corners=4pt,minimum width=2.8cm,minimum height=0.84cm,align=center,line width=0.92pt},
  io/.style={draw=themegray,fill=themegray!9!themebg,rounded corners=4pt,minimum width=2.55cm,minimum height=0.82cm,align=center,line width=0.9pt},
  flow/.style={->,draw=themefg!90,line width=1.05pt},
  support/.style={->,draw=themegray,dashed,line width=0.9pt},
  control/.style={->,draw=themealert,dashed,line width=0.95pt},
  persist/.style={->,draw=themegreen,dashed,line width=0.95pt}
]

\begin{pgfonlayer}{background}
  \fill[fill=themecurve!4!themebg,draw=themecurve!50,dashed,rounded corners=8pt]
    (1.7,-3.15) rectangle (13.95,5.2);
\end{pgfonlayer}
\node[anchor=west,font=\bfseries,color=themecurve] at (1.95,4.82) {Agent Harness / Runtime Contract};

% Main invocation path: the model is only one stage inside the harness.
\node[io]      (input)      at (0.15,1.25) {用户任务 / Event};
\node[runtime] (context)    at (4.0,1.25) {Context Builder\\\scriptsize 选择并组装当前工作集};
\node[model]   (model)      at (7.75,1.25) {LLM / Model\\\scriptsize Context $\rightarrow$ candidate output};
\node[guard]   (controller) at (11.45,1.25) {Runtime Controller\\\scriptsize validate · route · stop · execute};
\node[io]      (output)     at (15.55,1.25) {Approved Answer / Action};

\draw[flow] (input.east) -- (context.west);
\draw[flow] (context.east) -- (model.west);
\draw[flow] (model.east) -- node[above=1pt,font=\scriptsize]{候选输出 / Tool Call} (controller.west);
\draw[flow] (controller.east) -- (output.west);

% Control-plane inputs. They support the main path but do not become extra model capabilities.
\node[runtime] (registry) at (4.0,3.55) {Tool Registry\\\scriptsize schemas · implementations};
\node[runtime] (planner)  at (7.75,3.55) {Loop / Planner\\\scriptsize next call · stop condition};
\node[guard]   (policy)   at (11.45,3.55) {Policy / Permissions\\\scriptsize allow · deny · approve};

\draw[support] (registry.south) -- node[right=2pt,font=\scriptsize,color=themegray]{tool schemas} (context.north);
\draw[support] (registry.north) -- ++(0,0.72) -| node[pos=0.62,above=1pt,font=\scriptsize,color=themegray]{dispatch} (controller.north west);
\draw[control] (planner.south) -- ++(0,-0.55) -| (controller.north);
\draw[control] (policy.south) -- (controller.north east);

% Persistence plane. Durable state feeds context; committed effects and traces return here.
\node[state] (session) at (4.0,-1.15) {Session / State\\\scriptsize current task · checkpoint};
\node[state] (memory)  at (7.75,-1.15) {Memory / Artifacts\\\scriptsize persist · retrieve · reference};
\node[guard] (trace)   at (11.45,-1.15) {Trace / Eval Hooks\\\scriptsize calls · errors · latency · cost};

\draw[persist] (session.north) -- ++(0,0.55) -| (context.south west);
\draw[persist] (memory.north) -- ++(0,0.55) -| (context.south east);
\draw[persist] (controller.south) -- ++(0,-0.55) -| node[pos=0.72,below=1pt,font=\scriptsize,color=themegreen]{commit state} (session.east);
\draw[control] (controller.south) -- (trace.north);

\node[anchor=west,font=\scriptsize,color=themegray] at (1.95,-2.45)
  {主线只表示一次受控调用：Context Builder 组装输入，Model 产生候选，Runtime Controller 决定候选能否成为真实输出或动作。};
\node[anchor=west,font=\scriptsize,color=themegray] at (1.95,-2.88)
  {上方是控制面，下方是持久化与观测面；它们围绕模型工作，但不等同于模型本身。};

\end{tikzpicture}

\diagramcaption{Model 与 Agent Harness 的责任边界。主调用链是“Context Builder $\rightarrow$ Model $\rightarrow$ Runtime Controller”；上方控制面提供 Tool Registry、Loop/Planner 与 Policy，下方持久化面提供 Session/State、Memory/Artifacts 与 Trace/Eval。模型只产生候选输出，真实动作必须经过 Runtime Controller。}
\end{handbookdiagram}

一个 Harness 是否值得加入某个组件，要由真实 Failure 决定。例如：

\begin{itemize}
  \item 长任务跨 Context 丢失进度 $\Rightarrow$ 外部 artifact / checkpoint / compaction；
  \item 工具输出淹没主上下文 $\Rightarrow$ filesystem offload 或 subagent context isolation；
  \item 模型过早宣布完成 $\Rightarrow$ explicit completion criteria + verification；
  \item 写操作重复执行 $\Rightarrow$ idempotency key / checkpoint / side-effect ledger；
  \item 高风险动作越权 $\Rightarrow$ permission policy + human approval；
  \item 出错后无法解释 $\Rightarrow$ trace + structured state transitions。
\end{itemize}

因此 Harness 不是“功能越多越高级”。模型能力、任务结构和运行环境改变后，原来必要的 planning、memory、subagent 或 prompt scaffolding 可能变成多余成本。维护 Harness 时要持续重新验证每个组件是否仍然承担真实责任。

\subsection{长时程 Agent 的关键分离：Session 不等于 Context Window}

短任务里，conversation history、当前 Context 与“这次任务发生过什么”看起来几乎是一回事；一旦任务跨越多个 context windows、进程重启或多小时执行，这三个对象就必须拆开。

更稳定的模型是：

$$
\boxed{
\text{Session}
=\text{可恢复的事件历史与任务事实源}
}
$$

而第 $t$ 轮真正送入模型的只是它的一张工作视图：

$$
\boxed{
Context_t
=\operatorname{Assemble}
\left(
\operatorname{Select}(Session, State, Artifacts),
Memory,
Retrieved\ Evidence
\right)
}
$$

这里 \code{Select} 可以包含 slicing、trimming、compaction、summary、prompt-cache-friendly organization 等策略。因此：

\begin{itemize}
  \item \textbf{Session 要可恢复}：原始 events、tool results、状态变化和关键 artifacts 不应因为当前 Context 被压缩就永久消失；
  \item \textbf{Context 可以有损}：为了 token budget，可以只选择当前步骤真正需要的局部历史；
  \item \textbf{Compaction 不是持久化}：summary 只是新的 Context 表示，不能替代原始可审计历史；
  \item \textbf{Checkpoint 也不是完整 Session}：checkpoint 负责恢复当前执行状态，session log 负责保留更完整的事件轨迹。
\end{itemize}

Anthropic 在 2026 Managed Agents 的工程设计中把这层边界明确成 append-only session log：Harness 可以按位置重新读取历史 events，再选择怎样组织成新的 Context。这个实现不是唯一标准，但它揭示了一个可迁移原则：\textbf{可恢复历史应活在 Context Window 外部，Context Engineering 只负责构造当前工作集。}

\subsection{Brain--Hands--Session：长时程 Agent 的故障域要解耦}

把 Harness、执行环境和会话历史全部放进同一个容器，短期实现很直接，但会把三个本来不同的 failure domain 绑在一起。更稳的接口化结构是：

\begin{itemize}
  \item \textbf{Brain}：Model + Harness，负责读 Context、做决策、路由 Tool、处理 Observation；
  \item \textbf{Hands}：Sandbox / Workspace / MCP / External Tools，负责真实文件、代码、网络和业务副作用；
  \item \textbf{Session}：独立持久的事件日志与任务状态，用于恢复、回放和重新构造 Context。
\end{itemize}

\begin{handbookdiagram}
\centering
\begin{tikzpicture}[
  >=Stealth,
  font=\small,
  color=themefg,
  text=themefg,
  session/.style={draw=themecurve,fill=themecurve!12!themebg,rounded corners=5pt,minimum width=3.15cm,minimum height=0.92cm,align=center,line width=1pt},
  brain/.style={draw=themeamber,fill=themeamber!10!themebg,rounded corners=5pt,minimum width=3.15cm,minimum height=0.92cm,align=center,line width=1pt},
  hand/.style={draw=themegreen,fill=themegreen!10!themebg,rounded corners=5pt,minimum width=3.15cm,minimum height=0.92cm,align=center,line width=1pt},
  guard/.style={draw=themealert,fill=themealert!9!themebg,rounded corners=5pt,minimum width=3.15cm,minimum height=0.86cm,align=center,line width=0.95pt},
  flow/.style={->,draw=themefg!88,line width=1.02pt},
  durable/.style={->,draw=themecurve,line width=0.98pt,dashed},
  action/.style={->,draw=themegreen,line width=1.0pt},
  observe/.style={->,draw=themeamber,line width=0.98pt,dashed},
  policy/.style={->,draw=themealert,line width=0.95pt,dashed}
]

% Three failure domains. Horizontal position encodes ownership, not time.
\begin{pgfonlayer}{background}
  \fill[fill=themecurve!4!themebg,draw=themecurve!45,dashed,rounded corners=8pt]
    (-0.15,-1.25) rectangle (4.35,5.65);
  \fill[fill=themeamber!4!themebg,draw=themeamber!45,dashed,rounded corners=8pt]
    (4.75,-1.25) rectangle (10.85,5.65);
  \fill[fill=themegreen!4!themebg,draw=themegreen!45,dashed,rounded corners=8pt]
    (11.25,-1.25) rectangle (16.15,5.65);
\end{pgfonlayer}

\node[anchor=west,font=\bfseries,color=themecurve] at (0.1,5.25) {Durable Session};
\node[anchor=west,font=\bfseries,color=themeamber] at (5.0,5.25) {Brain = Model + Harness};
\node[anchor=west,font=\bfseries,color=themegreen] at (11.5,5.25) {Hands = Execution Environments};

% Durable session stores the history that survives any one context window.
\node[session] (log) at (2.1,3.75) {Append-only Event Log\\\scriptsize user · model · tool · state · error};
\node[session] (artifact) at (2.1,1.55) {Artifacts / Checkpoints\\\scriptsize spec · plan · files · committed progress};
\node[anchor=west,font=\scriptsize,color=themegray] at (0.2,-0.65)
  {可恢复事实源：Context 可以压缩或丢弃，Session 历史与关键产物仍保留。};

% Brain main path: select a view, decide, then govern the candidate action.
\node[brain] (context) at (6.25,3.75) {Context View$_t$\\\scriptsize select · slice · compact};
\node[brain] (model)   at (8.15,2.15) {LLM + Agent Loop\\\scriptsize decide next action / stop};
\node[brain] (runtime) at (9.15,0.35) {Harness Runtime\\\scriptsize validate · route · retry · resume};

\draw[durable] (log.east) -- node[above=1pt,font=\scriptsize,color=themecurve]{replay / select} (context.west);
\draw[durable] (artifact.east) -- ++(0.75,0) |- node[pos=0.7,below=1pt,font=\scriptsize,color=themecurve]{resume / reference} (context.south west);
\draw[flow] (context.south east) -- (model.north west);
\draw[flow] (model.south east) -- (runtime.north west);

% Hands are reached through one explicit execution gateway; generated code never owns credentials.
\node[hand] (gateway) at (13.65,3.55) {Execution Gateway\\\scriptsize provision · dispatch · collect result};
\node[hand] (sandbox) at (12.45,1.35) {Sandbox / Workspace\\\scriptsize files · shell · code};
\node[hand] (tools)   at (14.85,1.35) {External Tools / MCP\\\scriptsize DB · SaaS · browser · APIs};
\node[guard] (vault)  at (13.65,-0.45) {Credential / Egress Boundary\\\scriptsize scoped secrets · network policy};

\draw[action] (runtime.east) to[bend left=12] node[above=1pt,font=\scriptsize,color=themegreen]{execute(name,input)} (gateway.west);
\draw[observe] (gateway.west) to[bend left=12] node[below=1pt,font=\scriptsize,color=themeamber]{observation / error} (runtime.east);
\draw[action] (gateway.south west) -- (sandbox.north);
\draw[action] (gateway.south east) -- (tools.north);
\draw[policy] (vault.north west) -- (sandbox.south);
\draw[policy] (vault.north east) -- (tools.south);

% Every committed step returns to durable state before the next context view is assembled.
\draw[durable] (runtime.south) -- ++(0,-1.05) -- (0.45,-0.7) |- node[pos=0.79,left=2pt,font=\scriptsize,color=themecurve]{append event} (log.west);
\draw[durable] (runtime.south) -- ++(0,-0.55) -| node[pos=0.62,above=1pt,font=\scriptsize,color=themecurve]{checkpoint / artifact} (artifact.south);

\node[anchor=west,font=\scriptsize,color=themegray] at (4.95,-1.85)
  {主循环是 Session $\rightarrow$ Context View $\rightarrow$ Decision $\rightarrow$ Execution $\rightarrow$ Session；Tool result 不绕过 Context Assembly 直接“喂回”模型。};

\end{tikzpicture}

\diagramcaption{长时程 Agent 的 Brain--Hands--Session 解耦。Session 保存可恢复历史，Harness 从中选择当前 Context；Sandbox / Tools 作为可替换执行环境承接副作用。这样 Harness crash、sandbox crash 与 context compaction 不再等价于任务历史丢失。}
\end{handbookdiagram}

这一分离带来三个重要结果：

\begin{enumerate}
  \item \textbf{Recoverability}：Harness 进程可以无状态重启，再从 Session 恢复；Sandbox 失败可以作为一次 Tool Error 返回，并重新 provision；
  \item \textbf{Security}：credentials 可以留在 sandbox 外部的 vault / proxy，由 Runtime 代理调用，而不是暴露给 Agent 生成的任意代码；
  \item \textbf{Scalability}：一个 Brain 可以按需连接多个 Hands，执行环境也不必在任务开始前全部预热。
\end{enumerate}

这里最值得记住的不是某个 hosted Agent API，而是接口原则：

$$
\boxed{
\text{durable state}
\;|\;
\text{decision process}
\;|\;
\text{side-effect environment}
}
$$

这里的竖线只表示\textbf{生命周期与 failure domain 应显式分离}，不表示统计独立或数学正交。三者可以合作，但不应因为实现方便就被迫共享同一个生命周期。

\subsection{Harness Scaffolding 是关于“模型目前做不到什么”的可检验假设}

Harness 中的每个 workaround 都隐含一个模型假设。例如“模型接近 context limit 会过早收尾”可能促使系统定期 context reset；“模型不会可靠检查自己的实现”可能促使系统加入独立 evaluator；“模型不知道如何拆长任务”可能促使系统强制 planner / todo。

这些假设不能永久冻结。Anthropic 的长任务实验提供了一个很好的例子：早期对 Sonnet 4.5 有效的 context reset，在后续更强模型上变成额外的 orchestration / latency 开销，因此被移除并改由自动 compaction 支撑连续 session。这个结果不能推广成“以后都不需要 reset”，但它说明：

\begin{criterion}{Harness 组件的保留条件}
每次更换 Model、Tooling、Context Window、Sandbox 或任务分布后，都应该在固定 eval set 上重新做 ablation：移除某个 Harness 组件后，Task Success、Safety、Cost、Latency 与 Recovery 是否真的变差？若没有，就应考虑删除这层 scaffolding。
\end{criterion}

因此现代 Agent Engineering 不是不断向 Harness 加功能，而是维持一组\textbf{最小、可解释、能被失败证据证明有必要}的运行时约束。

\section{Skill、Tool、MCP 与 Subagent：四种接口不要混成“Agent 能力”}

它们经常一起出现，但回答的是不同问题：

\begin{handbooktable}{L{0.16\linewidth}L{0.29\linewidth}Y}
\toprule
\textbf{对象} & \textbf{核心问题} & \textbf{稳定边界}\\
\midrule
Skill & 这类任务\textbf{应该怎样做}？ & 可复用的程序性知识、说明、模板、脚本与资源；本身不等于一次外部副作用\\
Tool & 当前这一步\textbf{可以执行什么动作}？ & 带 schema / contract 的可调用能力；真正副作用由 Runtime / Tool implementation 执行\\
MCP & 外部能力与上下文\textbf{怎样以标准协议暴露}？ & 标准化 Tools / Resources / Prompts 等协议接口；不决定任务规划、授权策略或是否应该调用\\
Subagent & 哪个子任务\textbf{值得隔离到另一个执行上下文}？ & 一个被委派明确目标的独立 agent run，可有自己的 prompt、tools、context 和 budget；不要求一定并行\\
\bottomrule
\end{handbooktable}

把 Skill 类比为 procedural memory 有教学价值，因为它保存“以后遇到这类任务怎样行动”的稳定知识；但二者不能直接画等号。Skill 通常是显式、可版本控制的工程资产，还可以打包脚本与参考资料；“procedural memory”则是更宽的认知/系统组织概念。

以当前 Anthropic Agent Skills 为代表实现，一个 Skill 目录以 \code{SKILL.md} 为入口：name / description 等 metadata 先帮助 Agent 判断相关性，正文在任务匹配时再读取，更深的 reference / script / resource 继续按需加载。这种 progressive disclosure 的价值不是形式漂亮，而是避免把所有专业知识永久塞进 Context。

Tool 与 Skill 也不能互换。一个“部署 FastAPI 服务”的 Skill 可以告诉 Agent 应检查哪些文件、依次执行哪些步骤、怎样验证结果；真正的 \code{bash}、GitHub、Cloud API 或文件写入仍然是 Tools。Skill 教流程，Tool 提供动作。

Subagent 则解决另一类问题：当一个子任务会产生大量中间搜索、需要专门工具/指令，或应与主任务隔离 Context 时，把它委派给独立 worker，主 Agent 只接收压缩结果或 artifact reference。简单单步任务不应为了“多智能体”而创建 subagent。

\subsection{Defaults、Examples 与 Feedback：把领域经验变成可维护行为}

Agent 的行为质量不能只靠让最终用户不断“写更好的 Prompt”。对高频、可重复的领域任务，系统设计者通常应该先提供\textbf{经过验证的默认行为}，再开放必要的配置面。这背后的原因不是用户不会写 Prompt，而是系统 Owner 才拥有完整的业务约束、风险边界、历史失败案例和评测集。

稳定的领域行为可以分成三层：

\begin{itemize}
  \item \textbf{Defaults / Policy}：默认模型、工具权限、回答格式、必须检查的字段、停止条件和升级人工的规则；
  \item \textbf{Examples}：代表真实分布的正例、反例和边界例，用来说明“在这种输入下什么结果才算对”；
  \item \textbf{Feedback}：人工修正、失败 case、用户拒绝或业务结果，用来发现现有规则哪里失效。
\end{itemize}

三者的生命周期不能混在一起。一次人工修正只是 evidence，不应直接自动改写生产 Prompt、Skill 或 Long-term Memory。更稳妥的路径是：

\begin{processchain}
  \processstage{Observed Failure / Human Correction}
  \processstage{归因到具体失败层}
  \processstage{加入 Regression Case}
  \processstage{修改 Prompt / Skill / Tool / Workflow}
  \processstage{重跑 Eval}
  \processstage{通过后发布}
\end{processchain}

这样“系统会学习”才有可审计含义：它不是看到一次反馈就永久改变行为，而是把反馈转换成可复现的测试证据，再决定是否升级 Canonical 行为。对于安全敏感、高价值业务尤其如此。

Examples 也不是越多越好。重复、互相矛盾或只覆盖 happy path 的示例会增加 Context 成本并制造错误先验。选择示例时应覆盖关键分支和失败边界，并用真实 eval 证明加入这些 examples 后整体任务成功率提高，而不是只挑一个演示案例。

\section{State 与 Memory：先区分“当前执行”与“以后还能记得”}

一个独立的语言模型调用只根据这一次请求中可见的 Context 生成输出；应用跨轮次“记住”某件事，通常不是因为模型权重在每次聊天后被更新，而是因为运行时把信息持久化，并在后续调用中重新取回。某些托管 API 可以替应用保存 conversation/thread state，但这仍属于运行时状态管理，不等于模型参数获得了新的长期记忆。

先把五个容易混淆的对象分开：

\begin{handbooktable}{L{0.25\linewidth}Y}
\toprule
\textbf{对象} & \textbf{真正保存什么}\\
\midrule
Conversation History & 某段对话中的消息、模型输出与工具交互记录\\
Graph / Task State & 当前 workflow 为继续执行而需要的结构化状态\\
Checkpoint & 某个执行步骤的 State Snapshot，用于恢复、回放或人工中断\\
Long-term Memory & 跨 thread / session 仍希望被后续任务重新调用的信息\\
Vector Store & 一种可检索索引/存储方式；可以服务文档 RAG，也可以服务 Memory，但本身不等于 Memory\\
\bottomrule
\end{handbooktable}

LangChain v1 Agent 的 short-term memory 以 Agent State + Checkpointer 为主；跨 thread 的持久信息可以放入 Store。LangGraph persistence 通过 checkpoint 支持 human-in-the-loop、fault recovery、replay / time-travel debugging 和 thread-level memory。开发阶段可以使用 \code{InMemorySaver}；生产环境需要数据库支持的持久化实现。Checkpoint 能恢复 workflow state，但不能自动解决业务系统里的重复写入和外部副作用。

\subsection{Long-term Memory 的母模型：写入、维护、检索、注入是四个决策点}

把一切历史记录都保存下来，不等于建立了可用的记忆系统。一个真正的长期记忆闭环至少包含：

\begin{processchain}
  \processstage{候选证据}
  \processstage{Write Policy}
  \processstage{Durable Store}
  \processstage{Maintenance}
  \processstage{Retrieval Policy}
  \processstage{Context Injection}
\end{processchain}

这里的箭头表示数据进入下一阶段的运行顺序；每一阶段都可以由确定性代码、模型判断或人工控制，并不要求全部交给同一个 Agent。

\begin{handbookdiagram}
\centering
\begin{tikzpicture}[
  >=Stealth,
  font=\small,
  color=themefg,
  text=themefg,
  source/.style={draw=themegray,fill=themegray!10!themebg,rounded corners=4pt,minimum width=2.85cm,minimum height=0.9cm,align=center,line width=0.92pt},
  gate/.style={draw=themealert,fill=themealert!10!themebg,rounded corners=4pt,minimum width=2.85cm,minimum height=0.9cm,align=center,line width=0.98pt},
  store/.style={draw=themecurve,fill=themecurve!12!themebg,rounded corners=5pt,minimum width=3.55cm,minimum height=1.2cm,align=center,line width=1.05pt},
  maintain/.style={draw=themeamber,fill=themeamber!11!themebg,rounded corners=4pt,minimum width=3.15cm,minimum height=0.95cm,align=center,line width=0.98pt},
  context/.style={draw=themegreen,fill=themegreen!11!themebg,rounded corners=4pt,minimum width=2.85cm,minimum height=0.9cm,align=center,line width=0.98pt},
  flow/.style={->,draw=themefg!90,line width=1.02pt},
  write/.style={->,draw=themealert,dashed,line width=0.98pt},
  read/.style={->,draw=themegreen,line width=0.98pt},
  maintenance/.style={->,draw=themeamber,dashed,line width=0.95pt},
  constraint/.style={->,draw=themegray,densely dashed,line width=0.9pt}
]

% Central store. Left is the write path; right is the read path; top is maintenance.
\node[source] (evidence) at (1.55,3.45) {Candidate Evidence\\\scriptsize interaction · tool result · event};
\node[gate]   (writegate) at (4.8,3.45) {Write Policy\\\scriptsize should persist? scope?};
\node[store]  (memory) at (8.55,3.45) {Durable Memory\\\scriptsize semantic · episodic · procedural};
\node[gate]   (readgate) at (12.45,3.45) {Retrieval Policy\\\scriptsize need memory? how to search?};
\node[context] (context) at (15.75,3.45) {Context Injection\\\scriptsize selected memory only};

\draw[flow] (evidence) -- (writegate);
\draw[write] (writegate) -- node[above=1pt,font=\scriptsize,color=themealert]{accepted write} (memory);
\draw[read] (memory) -- node[above=1pt,font=\scriptsize,color=themegreen]{query / candidates} (readgate);
\draw[read] (readgate) -- node[above=1pt,font=\scriptsize,color=themegreen]{selected records} (context);

% Retrieval is conditioned on the current task, not only on what is stored.
\node[source] (task) at (12.45,1.25) {Current Task / Turn\\\scriptsize query · intent · current state};
\draw[constraint] (task.north) -- node[right=2pt,font=\scriptsize,color=themegray]{触发 / 约束检索} (readgate.south);

% Maintenance reads the store, reasons about lifecycle, then commits an update.
\node[maintain] (maintain) at (8.55,6.0) {Maintenance / Reflection\\\scriptsize correct · supersede · expire · merge · delete};
\draw[maintenance] (memory.north east) -- ++(0.75,0.55) |- node[pos=0.72,right=2pt,font=\scriptsize,color=themeamber]{inspect} (maintain.east);
\draw[maintenance] (maintain.west) -- ++(-0.9,0) |- node[pos=0.62,left=2pt,font=\scriptsize,color=themeamber]{commit update} (memory.north west);

% Metadata constrains all lifecycle decisions; it is not another memory type.
\node[source,minimum width=3.8cm] (meta) at (3.65,6.0) {Scope + Provenance + Time\\\scriptsize user · project · source · validity};
\draw[constraint] (meta.south west) -- ++(0,-0.55) -| (writegate.north);
\draw[constraint] (meta.east) -- (maintain.west);
\draw[constraint] (meta.south east) -- ++(0,-1.05) -| (readgate.north west);

\node[anchor=west,font=\scriptsize,color=themegray] at (0.1,7.05)
  {长期记忆不是一条“保存历史”的流水线：写入、维护、检索、注入分别有自己的触发条件与失败模式。};
\node[anchor=west,font=\scriptsize,color=themegray] at (0.1,0.05)
  {Reject / Skip 没有画成进入 Memory 或 Context 的箭头；未通过 gate 的信息就停在该决策点。};

\end{tikzpicture}

\diagramcaption{Agent 长期记忆的三个生命周期接口。左侧写入路径决定哪些证据可以进入 Durable Memory；右侧读取路径由 Current Task 触发检索并只把选中的记录注入 Context；上方维护路径负责更新、失效、合并和删除。Scope、Provenance 与 Time 同时约束写入、维护和检索。}
\end{handbookdiagram}

这张图最重要的结论是：

$$
\boxed{
\text{Memory}
\neq
\text{Conversation Log}
\neq
\text{Retrieval Index}
}
$$

Conversation log 可以是生成记忆的原始证据，vector / FTS / graph 可以是检索记忆的实现，而 Memory 是应用决定长期保留并在未来任务中重新使用的状态。

\subsection{第一维：记什么——Semantic、Episodic、Procedural}

Agent 工程里常借用三类记忆名称来组织不同内容。它们是实用的工程分类，不应被误写成对人类认知记忆系统的完整模型。

\begin{handbooktable}{L{0.22\linewidth}L{0.28\linewidth}Y}
\toprule
\textbf{类型} & \textbf{主要保存什么} & \textbf{Agent 例子}\\
\midrule
Semantic Memory & 相对稳定的事实、偏好、关系与当前已知状态 & 用户偏好暗色模式；某项目使用 PostgreSQL；Alex 负责 ML team\\
Episodic Memory & 带时间和情境的过去经历、事件或成功/失败轨迹 & 上周一次部署失败的上下文与处理结果；某次对话中已经完成的动作\\
Procedural Memory & “应该怎样做”的规则、技能、提示与操作流程 & 系统指令、SKILL.md、固定 SOP、从成功轨迹提炼出的执行步骤\\
\bottomrule
\end{handbooktable}

这三个类别回答的是“\textbf{记忆内容的语义是什么}”，不是“用什么数据库存”。同一个 Semantic Memory 可以放在 Markdown、SQLite、Postgres、vector index 或 graph 中；同一个数据库也可以同时存 semantic 与 episodic records。

\subsection{第二维：怎么存——五种常见架构模式可以叠加}

实际系统里常见下面五种模式。它们并不完全处在同一个抽象层：前三种更接近存储/检索表示，后两种把抽取、更新或关系建模提升成更完整的 memory layer。因此不要把它们理解成互斥的五选一。

\begin{handbooktable}{L{0.20\linewidth}L{0.27\linewidth}Y}
\toprule
\textbf{模式} & \textbf{核心机制} & \textbf{适合与代价}\\
\midrule
Bounded Prompt / File Memory & 少量高价值事实保存在 Markdown / profile，并在会话开始时直接注入 Context & 最透明、易人工编辑；但受 Context/token 预算限制，过长后会增加干扰，且需要主动清理旧信息\\
Relational + Lexical Memory & SQLite/Postgres 行记录 + metadata/time + FTS/BM25 关键词检索 & 精确字段、日期、ID 和审计友好；对改写、同义词与跨语言表达的召回能力有限\\
Vector-backed Semantic Memory & 文本/记录生成 embedding，使用 cosine / dot-product 等相似度召回，并结合 metadata filter & 擅长语义近邻和改写；需要管理 embedding model、索引、阈值、误召回和内容更新后的向量同步\\
Managed Memory Layer & 由 memory manager 从交互中抽取、合并、更新、删除或检索记忆，底层可再接普通 Store / Vector Store & 少写业务胶水，但多了一层 LLM 判断、延迟、成本和错误写入风险；必须评估 extraction/update policy\\
Temporal Knowledge Graph & 把 entity / relation / episode 与有效时间组织成图，再结合 lexical/semantic search 与 graph traversal & 适合关系、多跳和随时间变化的事实；抽取、实体消歧、图维护和查询更复杂，小型偏好记忆可能明显过度设计\\
\bottomrule
\end{handbooktable}

当前系统可以帮助看清这些模式怎样组合：

\begin{itemize}
  \item \textbf{Hermes Agent} 当前把 \code{MEMORY.md} / \code{USER.md} 作为有界、人工可读的持久记忆，并在 session start 注入一个 frozen snapshot；完整历史另外保存在 \code{~/.hermes/state.db}，通过 SQLite FTS5 做 session search。也就是说“prompt memory”和“可搜索历史”是两条不同路径。
  \item \textbf{Waku Agent} 当前默认用 \code{.waku/state.db} 保存 facts / episodes，并用 FTS5 做关键词检索，同时生成可读的 \code{MEMORY.md} mirror；它把 semantic、episodic、procedural memory、retrieval gate 与 consolidation 明确拆成模块，并允许把 semantic store 升级到 Supabase pgvector、Mem0、Zep 或 LangMem 等后端。
  \item \textbf{Supabase + pgvector} 仍然是 Postgres：\code{pgvector} 给普通表增加 vector column、distance operator 与 ANN index。它可以承载 semantic memory，但“关系数据库”和“向量检索”并不冲突。
  \item \textbf{LangMem} 更像 memory-management toolkit：它区分 semantic / episodic / procedural memory，支持 hot-path 与 background formation，并可以通过 LangGraph BaseStore 连接持久化 Store。它不是一种固定数据库形态。
\end{itemize}

\subsection{第三维：怎么找——Preload、Lexical、Semantic、Graph/Hybrid}

Memory 是否能保存，和“这一轮能不能找到它”是两个问题。常见 retrieval path 可以分为：

\begin{enumerate}
  \item \textbf{Preload / Always-in-Context}：极小、极高价值的 profile / instruction 直接随每轮 Context 注入；没有搜索延迟，但会永久占用 token，并可能让无关旧记忆持续影响模型。
  \item \textbf{Lexical Retrieval}：SQLite FTS5、Postgres full-text、BM25 等根据 token/词项匹配；适合名字、ID、错误码、日期和精确关键词。
  \item \textbf{Semantic Retrieval}：将 query 与 memory record 映射为同一 embedding space，再做相似度检索；适合同义改写和语义近邻，但 embedding 相似并不等于事实相关或事实仍然有效。
  \item \textbf{Graph / Hybrid Retrieval}：用 entity、relation、时间和 graph distance 扩展或 rerank 候选，并可与 lexical / semantic retrieval 组合。图数据库本身并不要求“所有节点和边都必须做 embedding”；向量只是其中一种检索信号。
\end{enumerate}

因此一个很实用的设计是先加 \textbf{Retrieval Gate}：判断当前 turn 是否真的需要长期记忆，再决定是否支付搜索、embedding、rerank 和 Context token 成本。Waku 当前就把这个 gate 作为独立模块；但 gate 自身也有一次判断成本，而且 false negative 会导致“明明记得却没有取回”。

\subsection{第四维：怎么维护——不要把“新事实”简单等同于覆盖旧行}

长期运行以后，Memory 的主要难题通常从“存不下来”转成“旧事实、重复事实和冲突事实越来越多”。维护操作至少要区分：

\begin{itemize}
  \item \textbf{No-op / Reject}：这条信息没有长期价值，或来源不可信，不进入 memory；
  \item \textbf{Insert}：新增独立事实或 episode；
  \item \textbf{Update / Correct}：旧记录本身错误，需要改正；
  \item \textbf{Supersede / Invalidate}：旧事实曾经正确，但现在不再有效；保留历史，并记录新的有效区间；
  \item \textbf{Expire / TTL}：信息只在有限时间内有价值，到期后不再参与默认检索；
  \item \textbf{Delete}：因为用户请求、隐私、合规、错误数据或保留策略真正删除；
  \item \textbf{Consolidate / Reflect}：把多条重复、零碎记录合并成更稳定的事实、episode summary 或 procedural rule。
\end{itemize}

例如“产品发布日原定 5 月 20 日，后来改为 6 月 3 日”不是简单的真假冲突。若未来要回答“5 月初我们当时计划什么时候发布”，旧值仍有历史意义。Temporal memory 因此需要至少区分：

$$
\text{recorded/observed time}
\qquad\text{与}\qquad
\text{fact validity time}.
$$

Zep / Graphiti 的 temporal graph 就会在 relation/fact 上维护类似 \code{valid\_at} / \code{invalid\_at} 的时间字段，使新信息可以 supersede 旧关系而不必直接抹去历史。它的当前检索也可以结合 semantic、keyword 和 graph-distance signals；这比“图数据库就是把节点和边都嵌入以后做向量搜索”更准确。

\subsection{Hot Path 与 Background：记忆形成发生在什么时候}

写 Memory 也有运行时选择：

\begin{itemize}
  \item \textbf{Hot Path}：在当前交互中立即判断并写入。关键偏好或用户明确说“记住这个”时很合适；缺点是增加响应延迟，而且当前对话尚未结束时可能过早总结。
  \item \textbf{Background / Reflection}：对话结束或积累若干 turns 后再批量抽取、去重、合并与更新。不会阻塞前台回答，更容易获得完整上下文；代价是记忆有延迟，后台任务还要处理重试、重复执行和最终一致性。
\end{itemize}

LangMem 当前明确支持 hot-path memory tools 与 background \code{ReflectionExecutor} / memory manager；Waku 也把若干聊天后的 consolidation 与当前 turn retrieval 分开。所谓“睡眠/反思”最稳妥的工程含义就是：\textbf{把不需要阻塞用户的 memory maintenance 从主请求路径移到后台处理}，而不是假设 Agent 在离线时会自动获得更强的认知能力。

\subsection{当前实现差异：产品名不能替代架构判断}

版本变化很快，尤其不能把某个产品旧版的 memory policy 写成永久事实：

\begin{itemize}
  \item \textbf{Mem0 OSS v3（2026）} 当前已经改为 single-pass ADD-only extraction；检索融合 semantic、BM25 keyword 与 entity matching。旧版“同一抽取器在 ADD / UPDATE / DELETE 之间决策”的描述不再适合作为当前 OSS 默认行为。
  \item Mem0 的旧 OSS external graph-store integration 已移除；当前 Graph Memory 属于 Mem0 Platform 内建能力。因此“Mem0 OSS 默认同时维护外部 Neo4j 图”已经过时。
  \item \textbf{LangMem} 的 memory manager 仍可按配置执行 insert / update / delete，而且它把 memory type、memory formation 与底层 Store 分开；因此它和 Mem0 当前 OSS 的 write policy 不能混成同一个模型。
  \item \textbf{Zep / Graphiti} 的价值重点在 temporally-aware knowledge graph，而不是“图一定比向量更高级”。对于少量独立偏好，图抽取与异步 ingest 的复杂度可能没有收益；只有关系、多跳和时间演化确实重要时才值得承担这层成本。
\end{itemize}

单次 demo 中某个 backend 用了 4 秒、10 秒甚至更久，只能说明那个具体运行的 model、network、ingestion path 和数据规模；不能直接推成“某类 memory architecture 永远更慢”。要比较系统，至少固定 corpus、query set、LLM/embedding model、top-k、写入生命周期和评测方法，并报告多次运行的 latency distribution。

\subsection{Memory 的安全边界：持久化错误比一次 hallucination 更危险}

Memory 会在未来反复进入 Context，所以错误写入、恶意写入和跨用户泄漏具有\textbf{持久化放大效应}。生产系统至少要处理：

\begin{itemize}
  \item \textbf{Scope isolation}：user / tenant / project / agent namespace 必须明确，不能让一个用户的 memory 被另一个用户检索到；
  \item \textbf{Provenance}：记录信息来自用户明确陈述、工具结果、业务数据库、网页还是模型推断，必要时保留 source id / timestamp；
  \item \textbf{Trust boundary}：网页、邮件、工具返回中的 prompt injection 不能未经审查直接升级为 procedural memory 或高优先级 instruction；
  \item \textbf{Sensitive data policy}：凭据、密钥、受限 PII 不应因为“以后可能有用”就进入普通 memory store；
  \item \textbf{Deletion / retention}：用户撤回、法规要求、项目结束和错误修正必须有真实删除/失效路径；
  \item \textbf{Auditability}：应能回答谁写了这条 memory、何时写、后来怎样被更新、哪一次回答检索并使用了它。
\end{itemize}

对于会直接进入 system prompt 的 curated memory，write gate 应比普通文档索引更严格，因为一条恶意或错误 instruction 可能跨很多 future turns 持续影响行为。

\subsection{Memory 怎么评：不能只问“最终回答看起来对不对”}

至少把评测拆成五层：

\begin{handbooktable}{L{0.24\linewidth}Y}
\toprule
\textbf{评测层} & \textbf{要回答的问题}\\
\midrule
Write Quality & 应该记的是否被捕获；不该记的是否被拒绝；抽取事实是否忠实于来源\\
Maintenance Quality & 更新、supersede、expire、delete、deduplicate 是否处理正确；旧事实会不会错误压过新事实\\
Retrieval Quality & 需要时能否召回；top-k 中无关 memory 有多少；跨语言/改写/时间查询是否稳定\\
End-to-end Utility & 注入 memory 后任务成功率是否提高；是否反而引入偏见、过度解释或 stale-memory error\\
System Cost / Safety & write/read latency、额外 tokens、embedding/LLM cost、跨租户泄漏、memory poisoning 与删除合规是否可接受\\
\bottomrule
\end{handbooktable}

因此“记得更多”不是目标。更好的目标是：\textbf{只把未来任务真正需要、来源可追溯、生命周期清楚的信息，在需要的时候准确取回，并以可控成本注入 Context。}

\section{Loop 与 Graph：开放探索和显式流程不是二选一}

“Loop engineering”和“Graph engineering”容易被讲成两代互相替代的 Agent 技术，实际更稳定的区别是：\textbf{下一步控制权在运行时由模型决定多少，又有多少已经被系统设计者显式写进控制流。}

\begin{handbooktable}{L{0.20\linewidth}L{0.31\linewidth}Y}
\toprule
\textbf{结构} & \textbf{控制方式} & \textbf{适合的问题}\\
\midrule
Agent Loop & Model 根据当前 State / Observation 反复选择下一动作，直到满足 stop condition & 开放搜索、调研、调试、未知工具序列；下一步无法预先稳定确定\\
Workflow / Graph & 主要阶段、节点、边、并行关系和部分分支在设计时显式定义 & 审批、退款、ETL、固定业务流程；需要可预测路径、恢复和审计\\
Hybrid & Graph 规定大结构，其中某些 Node 内运行 Agent Loop；Loop 也可调用一个固定 subworkflow & 生产系统最常见：稳定业务骨架 + 少量开放问题求解\\
\bottomrule
\end{handbooktable}

Graph 也不等于“全部确定性”。Node 可以是普通函数、LLM call、Tool、router 或完整 Agent；Edge 也可以根据运行时 state 条件选择。Graph 真正提供的是\textbf{显式拓扑和状态转移边界}，不是消灭模型随机性。

\begin{handbookdiagram}
\centering
\begin{tikzpicture}[
  >=Stealth,
  font=\small,
  color=themefg,
  text=themefg,
  nodebox/.style={draw=themecurve,fill=themecurve!12!themebg,rounded corners=4pt,minimum width=2.25cm,minimum height=0.82cm,align=center,line width=0.95pt},
  tool/.style={draw=themegreen,fill=themegreen!10!themebg,rounded corners=4pt,minimum width=2.25cm,minimum height=0.82cm,align=center,line width=0.9pt},
  decision/.style={draw=themeamber,fill=themeamber!11!themebg,rounded corners=4pt,minimum width=2.15cm,minimum height=0.82cm,align=center,line width=0.95pt},
  embedded/.style={draw=themealert,fill=themealert!6!themebg,dashed,rounded corners=6pt,minimum width=2.7cm,minimum height=1.22cm,align=center,line width=0.95pt},
  flow/.style={->,draw=themefg!90,line width=1.0pt},
  dynamic/.style={->,draw=themeamber,line width=1.0pt,dashed}
]

% Left panel: an open loop. The model decides the next action after each observation.
\node[anchor=west,font=\bfseries,color=themeamber] at (0.1,5.05) {Agent Loop：下一步由运行时观察决定};
\node[nodebox] (lm) at (2.45,3.55) {Model Decision};
\node[tool] (action) at (2.45,1.95) {Tool / Action};
\node[nodebox] (obs) at (2.45,0.35) {Observation / State};
\node[decision] (done) at (5.05,3.55) {Stop Condition};

\draw[flow] (lm) -- node[right=2pt,font=\scriptsize]{选择动作} (action);
\draw[flow] (action) -- (obs);
\draw[dynamic] (obs.west) -- ++(-1.15,0) |- node[pos=0.62,left=2pt,font=\scriptsize,color=themeamber]{下一轮} (lm.west);
\draw[dynamic] (lm.east) -- (done.west);
\draw[flow] (done.east) -- ++(0.85,0) node[right,font=\scriptsize]{满足：结束};

% Right panel: the graph fixes the large-scale topology; one node may itself contain a loop.
\node[anchor=west,font=\bfseries,color=themecurve] at (7.45,5.05) {Graph / Workflow：大结构显式，局部仍可动态};
\node[nodebox] (start) at (8.25,3.35) {START};
\node[nodebox] (retrieve) at (10.65,3.35) {Retrieve};
\node[embedded] (agent) at (13.05,3.35) {Agent Loop Node\\\scriptsize decide $\rightarrow$ tool $\rightarrow$ observe $\looparrowleft$};
\node[nodebox] (verify) at (15.45,3.35) {Verify};
\node[nodebox] (end) at (15.45,1.15) {END};

\draw[flow] (start) -- (retrieve);
\draw[flow] (retrieve) -- node[above=1pt,font=\scriptsize]{开放子任务} (agent);
\draw[flow] (agent) -- node[above=1pt,font=\scriptsize]{artifact / result} (verify);
\draw[flow] (verify) -- node[right=2pt,font=\scriptsize]{通过} (end);

\node[anchor=west,font=\scriptsize,color=themegray] at (0.1,-0.75)
  {左图的循环边本身就是主控制结构；右图的 START / Retrieve / Agent Loop / Verify / END 已预先固定，只有 Agent Loop Node 内部保留开放决策。};
\node[anchor=west,font=\scriptsize,color=themegray] at (0.1,-1.18)
  {因此 Graph 与 Loop 的差异是控制权落在哪里，不是“确定性系统”和“智能系统”的二选一。};

\end{tikzpicture}

\diagramcaption{Loop 与 Graph 的控制权差异。左侧 Loop 的下一动作在运行时动态决定；右侧 Graph 先固定主要节点与边，但其中一个 Node 仍然可以包含完整 Agent Loop。箭头表示控制流，不表示理论上的技术演进顺序。}
\end{handbookdiagram}

因此第一判断不是“Loop 还是 Graph 哪个先进”，而是：

\begin{enumerate}
  \item 哪些步骤的顺序和合法分支已经由业务规则确定？先写进 Workflow / Graph；
  \item 哪些子问题必须根据未知 Observation 临场选择下一动作？把局部控制权交给 Agent Loop；
  \item 哪些写操作必须经过审批、幂等和固定补偿路径？不要让自由 Loop 拥有无限制副作用；
  \item 每个 Loop 的 stop condition、iteration / time / cost budget 是什么？
\end{enumerate}

这也是为什么“固定退货规则”通常不需要 Planner Agent，而 Deep Research、未知 bug 调试或开放网页调查更适合保留 Loop。

\section{Multi-Agent：只有任务真的可分解时，并行委派才有价值}

Subagent 只是一次隔离委派；Multi-Agent 则要求多个 Agent runs 共同完成同一最终任务。最常见的稳定结构是 orchestrator--worker：

\begin{processchain}
  \processstage{User Goal}
  \processstage{Orchestrator 分解任务}
  \processstage{Workers 独立执行}
  \processstage{Artifacts / Findings}
  \processstage{Synthesis / Verification}
  \processstage{Final Output}
\end{processchain}

其中 Orchestrator Own 的不是“把所有细节再做一遍”，而是 decomposition、worker boundary、budget、progress、gap detection 与 synthesis。Worker 应收到明确的 objective、scope、允许的 tools、期望输出格式和 stop condition；否则多个 worker 很容易重复同一搜索、遗漏边界或产生互相矛盾的结论。

\subsection{什么时候值得使用多个 Agent}

Multi-Agent 的主要收益来自\textbf{可并行的独立子问题与 context isolation}，而不是“多几个模型投票就一定更聪明”。优先考虑它的场景包括：

\begin{itemize}
  \item 一个问题天然分成多个相对独立的研究方向，可以并行搜索；
  \item 每个分支会产生大量中间 Tool outputs，不希望污染主 Agent Context；
  \item 不同分支需要不同 tools、instructions 或 models；
  \item 单个 context window 无法经济地容纳全部探索过程；
  \item wall-clock latency 可以通过真实并行显著降低，而且任务价值足以覆盖额外成本。
\end{itemize}

反过来，如果任务强依赖共享中间状态、步骤高度串行、子问题之间不断互相修改，或者本来一次 Tool Call 就能完成，多 Agent 只会增加 token、同步、权限和失败面。

总成本至少会变成：

$$
C_{\text{multi}}
= C_{\text{orchestrator}}
+\sum_{i=1}^{k} C_{\text{worker},i}
+C_{\text{synthesis}}
+C_{\text{coordination}}.
$$

Anthropic 在其 2025 Research 系统的内部数据中报告：多 Agent Research 相比普通 chat 使用了约 $15\times$ tokens，并在特定 internal research eval 上获得明显收益；这只能作为“并行研究可能用成本换覆盖”的工程证据，不能推广成多 Agent 对任意任务都更好。

\subsection{Deep Research：一个可复用的 Integration 母例}

Deep Research 最值得保留的不是某个产品的 Agent 名称，而是这条可生成的结构：

$$
\boxed{
\text{Clarify / Scope}
\to\text{Research Brief}
\to\text{Orchestrate}
\to\text{Parallel Evidence Gathering}
\to\text{Synthesis}
\to\text{Citation / Verification}
}
$$

第一步先把问题边界和交付物写清；随后 Orchestrator 根据 brief 动态拆分研究方向；Workers 各自在隔离 Context 中搜索、读取和分析；完成后返回\textbf{可验证 findings、sources 或 artifact references}，而不是要求系统永久保存或暴露每个 Agent 的私有思维过程。最后由 synthesizer 解决重复、冲突、缺口和引用，再决定是否需要追加一轮研究。

这套结构同时展示了 Loop 与 Graph 的混合：整体生命周期可以显式编排，但“研究哪个来源、发现新线索后往哪里继续”仍属于 Worker Loop 的动态控制。

Multi-Agent 的正确性也必须分层验证：worker 是否完成自己边界内的任务；handoff 是否丢信息；synthesizer 是否忠实引用 worker artifacts；最终输出是否满足全局目标。只评最后一段文字，会看不见协调失败发生在哪一层。

\subsection{Planner--Generator--Evaluator：多 Agent 的另一条价值是分离“做”和“判”}

Multi-Agent 不只用于把搜索并行化。另一个高价值结构是把\textbf{生成工作}与\textbf{评价工作}放进不同 execution contexts：

\begin{processchain}
  \processstage{Planner 定义交付物}
  \processstage{Generator 实现一个可验证增量}
  \processstage{Evaluator 独立操作环境并检查}
  \processstage{Structured Feedback}
  \processstage{修复或进入下一增量}
\end{processchain}

这一结构针对的是 self-evaluation bias：同一个 Agent 在刚生成结果后再问自己“做得好吗”，很容易沿用已有解释并过度宽容。把 evaluator 隔离出来以后，可以给它独立 rubric、Tools 和失败阈值，让它从环境结果而不是生成者叙述出发判断。

对于软件工程，最强的 evaluator 往往不是“再读一遍代码”，而是运行 tests、浏览真实 UI、调用 API、检查数据库状态和重现用户路径。对于主观任务，也应尽量把“好不好”拆成可操作 criterion，例如 originality、craft、functionality，而不是只让另一个 LLM 给总分。

Anthropic 2026 的长时程应用实验进一步使用了 \textbf{sprint contract}：Generator 和 Evaluator 在写代码前先明确这一小段工作“完成”的可测试定义，再进入实现。它可以抽象成：

$$
\boxed{
\text{Goal}
\to\text{Acceptance Contract}
\to\text{Implementation}
\to\text{Independent Verification}
}
$$

这比“让多个 Agent 互相讨论直到满意”为何更稳：每一轮都有外部可检查的完成条件，feedback 能直接转成下一次修改，而不是继续扩散自然语言意见。

这种结构同样不是免费收益。独立 evaluator 会显著增加模型调用、工具执行和 wall-clock time；只有任务价值、失败成本或质量要求足以覆盖验证成本时才值得使用。

\section{LangChain 与 LangGraph：当前分工}

\subsection{LangChain v1：高层 Agent API}

当前 LangChain 的标准 Agent API 是 \codepath{langchain.agents.create_agent}：

\begin{lstlisting}[language=Python]
from langchain.agents import create_agent

agent = create_agent(
    model=model,
    tools=[search_docs, lookup_order],
    system_prompt="Use tools only when needed and cite evidence.",
)
\end{lstlisting}

\codepath{create_agent} 运行在 LangGraph runtime 上。LangChain 更适合快速组合标准 tool-calling agent、provider abstraction、middleware 与 structured output。

“代码少了”不代表系统复杂度消失。权限、状态、失败处理、评测和观测仍然由应用负责。

\subsection{Deep Agents：当前一个明确的 Harness 实现}

LangChain 当前把 \code{deepagents} 明确描述为 agent harness：底层仍是 tool-calling loop，但默认组合 planning、filesystem context management、subagent delegation、long-term memory、permission rules 与 human-in-the-loop 等能力，并运行在 LangGraph runtime 上。

这正好说明 Framework、Runtime 与 Harness 可以分层：LangGraph 提供 durable execution / state orchestration；LangChain 提供较高层 Agent API；Deep Agents 再提供更 opinionated 的 long-running Agent Harness。简单任务仍应优先使用更薄的 \codepath{create_agent} 或普通 LangGraph workflow，而不是默认加载全部 Harness 组件。

\subsection{LangGraph v1：显式 State + Node + Edge}

LangGraph 适合需要 multi-step workflow、conditional branch、loop、parallel branch、human approval、interrupt / resume、durable execution 和显式 state machine 的场景。当前官方文档也明确区分：workflow 的主要 code path 是预先定义的，而 agent 会在运行时动态决定 process / tool usage；两者都可以在同一个 LangGraph runtime 中组合。

核心模型可以写成：

$$
\boxed{\text{Node}: State_{in}\mapsto \Delta State}
$$

最小概念代码：

\begin{lstlisting}[language=Python]
from typing_extensions import TypedDict
from langgraph.graph import StateGraph, START, END

class State(TypedDict):
    query: str
    docs: list[str]

def retrieve(state: State):
    docs = search_index(state["query"])
    return {"docs": docs}

builder = StateGraph(State)
builder.add_node("retrieve", retrieve)
builder.add_edge(START, "retrieve")
builder.add_edge("retrieve", END)
graph = builder.compile()
\end{lstlisting}

标准分工可以压成：

\begin{itemize}
  \item High-level standard agent：LangChain \codepath{create_agent}；
  \item Custom stateful orchestration：LangGraph \code{StateGraph}。
\end{itemize}

如果任务只是确定性的“Input → Function → Output”，普通代码通常更简单，也更容易测试。

\section{MCP：标准化外部 Context 与 Capability 接口}

MCP（Model Context Protocol）解决的是：不同 LLM Host 怎样通过标准协议发现并调用外部 Context / Capability，而不是为每一个应用和每一个外部系统重新写一套私有集成。

最小角色是：

\begin{itemize}
  \item \textbf{Host}：运行 LLM 应用的宿主；
  \item \textbf{Client}：Host 内负责某条 MCP 连接的组件；
  \item \textbf{Server}：暴露 Context / Capability 的服务。
\end{itemize}

Server 的核心 primitives 仍包括 Tools、Resources 与 Prompts。

\begin{handbookdiagram}
\centering
\begin{tikzpicture}[
  >=Stealth,
  font=\small,
  color=themefg,
  text=themefg,
  hostnode/.style={draw=themecurve,fill=themecurve!12!themebg,rounded corners=4pt,minimum width=2.8cm,minimum height=0.9cm,align=center,line width=1pt},
  clientnode/.style={draw=themeamber,fill=themeamber!12!themebg,rounded corners=4pt,minimum width=2.55cm,minimum height=0.84cm,align=center,line width=0.95pt},
  servernode/.style={draw=themegreen,fill=themegreen!11!themebg,rounded corners=4pt,minimum width=3.05cm,minimum height=0.96cm,align=center,line width=1pt},
  backendnode/.style={draw=themegray,fill=themegray!10!themebg,rounded corners=4pt,minimum width=2.7cm,minimum height=0.84cm,align=center,line width=0.9pt},
  protocol/.style={<->,draw=themefg!88,line width=1pt},
  local/.style={->,draw=themegreen,line width=0.95pt,dashed},
  internal/.style={->,draw=themeamber,line width=0.95pt}
]

% Two responsibility regions. Position groups roles; it does not encode execution time.
\begin{pgfonlayer}{background}
  \fill[fill=themecurve!5!themebg,draw=themecurve!45,dashed,rounded corners=7pt]
    (-0.2,0.55) rectangle (6.15,8.3);
  \fill[fill=themegreen!5!themebg,draw=themegreen!45,dashed,rounded corners=7pt]
    (6.85,0.55) rectangle (15.6,8.3);
\end{pgfonlayer}

\node[anchor=west,font=\bfseries,color=themecurve] at (0.1,7.9) {Host Application};
\node[anchor=west,font=\bfseries,color=themegreen] at (7.15,7.9) {External MCP Servers};

% Host-side topology: application owns clients; each client owns one server connection.
\node[hostnode] (app) at (2.0,6.55) {LLM Application\\\scriptsize model / workflow};
\node[clientnode] (ca) at (4.3,5.15) {MCP Client A};
\node[clientnode] (cb) at (4.3,2.75) {MCP Client B};

% Separate internal corridors prevent Client A / Client B edges from crossing.
\draw[internal] (app.east) -- ++(0.55,0) -| (ca.north);
\draw[internal] (app.south) -- ++(0,-3.35) -- (cb.west);

% Client-server pairs are horizontally aligned; protocol edges are direct and unambiguous.
\node[servernode] (sa) at (9.25,5.15) {Server A\\\scriptsize Tools · Resources · Prompts};
\node[servernode] (sb) at (9.25,2.75) {Server B\\\scriptsize Tools · Resources · Prompts};
\draw[protocol] (ca.east) -- node[above=1pt,font=\scriptsize]{MCP} (sa.west);
\draw[protocol] (cb.east) -- node[above=1pt,font=\scriptsize]{MCP} (sb.west);

% Server-local backends are outside the MCP role model itself.
\node[backendnode] (dba) at (13.55,5.15) {Customer DB};
\node[backendnode] (svc) at (13.55,2.75) {Weather / SaaS API};
\draw[local] (sa) -- node[above=1pt,font=\scriptsize,color=themegreen]{本地集成} (dba);
\draw[local] (sb) -- node[above=1pt,font=\scriptsize,color=themegreen]{本地集成} (svc);

\node[anchor=west,font=\scriptsize,color=themegray] at (0.15,1.05)
  {这是角色与连接关系图，不是执行时序图：MCP 规范 Client ↔ Server 交互，Agent / Workflow 仍负责决定何时调用。};

\end{tikzpicture}

\diagramcaption{MCP 的静态角色边界。Host 内的 Client 与外部 Server 通过 MCP 协议交互；Server 再把自己的 Tool、Resource 或 Prompt 暴露给 Host。图中连接表示协议交互，不表示 Agent 的任务决策顺序。}
\end{handbookdiagram}

\subsection{当前 MCP 基线：2026-07-28}

截至 2026-08-22，当前 MCP 规范版本是 2026-07-28。新版核心采用 stateless request/response 模型：现代客户端通过 \code{server/discover} 获取服务器身份与能力，不再要求每个连接先建立 protocol-level session；Streamable HTTP 的现代路径也不再依赖 \code{Mcp-Session-Id} 绑定 worker。

新版还引入或强化了：

\begin{itemize}
  \item Multi Round-Trip Requests；
  \item header-based routing；
  \item cacheable list results；
  \item authorization hardening；
  \item Extensions framework；
  \item formal deprecation policy。
\end{itemize}

Roots、Sampling 与 MCP-level Logging 已进入 deprecated 路径；\code{ping} 被移出新协议。兼容旧 server 时，当前 Python SDK v2 的 Client 可以回退到 2025-era \code{initialize} handshake。

\subsection{MCP Python SDK v2}

当前官方 Python SDK 的稳定线是 v2，要求 Python 3.10+。高层 Server 类已经是 \code{MCPServer}：

\begin{lstlisting}[language=Python]
from mcp.server import MCPServer

mcp = MCPServer("Demo")

@mcp.tool()
def add(a: int, b: int) -> int:
    """Add two integers."""
    return a + b

if __name__ == "__main__":
    mcp.run()  # default transport: stdio
\end{lstlisting}

SDK 可以从函数名、docstring 与 type hints 推导 Tool 的名称、说明和 input schema。\code{@mcp.tool()} 的括号属于当前 API 契约。

当前主要 transport 包括：

\begin{itemize}
  \item \code{stdio}：本地 child process；
  \item Streamable HTTP：remote server；
  \item SSE：仍存在兼容支持，但已经是旧一代路径，不应作为新系统首选。
\end{itemize}

\code{stdio} 模式下 stdout 是协议通道，业务日志应走 stderr / logging，避免污染 wire protocol。

\subsection{Tool Schema 与错误边界}

Tool 至少需要结构化 input schema。当前 Python SDK 可以从 type hints 自动生成 JSON Schema。

可由模型调整参数后重试的业务错误，应该返回可操作的 execution feedback；协议结构本身错误则属于更低层的 RPC / protocol failure。两类错误如果混在一起，模型和 Runtime 都很难知道该重试、改参数还是直接终止。

\section{Security：任何 RAG / Tool / MCP 系统都要画 Trust Boundary}

真实 Agent 会同时接收用户输入、检索内容和 Tool / MCP 结果。这三类输入都可能不可信。

主要攻击面包括：

\begin{itemize}
  \item Prompt Injection；
  \item Retrieval Poisoning；
  \item Tool Injection；
  \item Data Exfiltration；
  \item Over-permissioned Tool；
  \item Secret / PII Leakage；
  \item Arbitrary Code Execution；
  \item Malicious MCP Server / Tool Result。
\end{itemize}

\subsection{先限制 Blast Radius，再讨论模型会不会犯错}

Agent 风险可以用一个非常粗但有用的分解来思考：

$$
\boxed{
Risk\approx P(\text{unsafe action})\times BlastRadius
}
$$

Prompt、Guardrail、模型训练与人工审批主要尝试降低前一项；Sandbox、filesystem boundary、network egress policy、least privilege、credential isolation 与 transaction boundary 主要限制后一项。因为任何概率型检测都可能漏掉，生产系统不能把“模型通常会拒绝”当成唯一安全边界。

这也解释了为什么高自主 Agent 越来越强调\textbf{environment-level containment}：

\begin{itemize}
  \item 不让不可信代码所在 Sandbox 看见长期 credentials；
  \item 文件系统只挂载当前任务真正需要的目录；
  \item 网络访问使用 allowlist / proxy / egress policy；
  \item 高风险 Tool 使用独立授权域和 side-effect ledger；
  \item Tool Result 仍按不可信外部输入检查，因为可信 Tool 也可能返回被注入的网页、README 或用户数据。
\end{itemize}

Human approval 很重要，但不应成为唯一防线。高频 permission prompt 会产生 approval fatigue；更稳的结构是把低风险动作放进受限环境自动执行，把真正高影响且用户能判断的少数决策保留给人工确认。

\begin{warning}{Trust Boundary｜进入 Context 的文本不自动拥有指令权限}
外部文档、网页、数据库字段和 Tool result 首先都是数据。系统必须把 Trusted Policy / Instructions 与 Untrusted Data / Evidence 分开，不能因为一段文本进入 Context 就默认它可以覆盖上层指令。
\end{warning}

生产 Tool 至少需要：input validation、authentication、authorization、least privilege、timeout、retry policy、rate limit、audit log、side-effect classification、idempotency strategy、rollback / compensation 与必要的 user confirmation。

\codepath{search_docs()} 这样的只读查询，与 \codepath{send_email()}、\codepath{charge_card()}、\codepath{delete_file()}、\codepath{update_database()} 这样的有副作用动作，不能共享同一授权和重试策略。

对于 HTTP-based MCP server，授权应遵循 MCP 当前 authorization contract；\code{stdio} transport 通常从运行环境获得 credentials，不走相同 HTTP OAuth 流程。无论使用哪种 transport，都必须把权限缩到完成任务所需的最小范围。

\section{Observability、Evaluation、Cost 与 Failure Recovery}

\subsection{Observability：必须能重建一次执行轨迹}

每次 Agent / Workflow 执行至少应该能够追踪：

\begin{itemize}
  \item User Request；
  \item Context Snapshot / References；
  \item Retrieved Chunks；
  \item Model Calls；
  \item Tool Calls / Arguments / Results；
  \item State Transitions；
  \item Retries / Errors；
  \item Token Usage；
  \item Latency / Cost；
  \item Final Output。
\end{itemize}

否则只能看到最后答案，无法判断错误发生在 Retrieval、Model、Tool、State 还是 Routing。

\subsection{Evaluation：按层拆开，而不是只看“回答准确率”}

\begin{handbooktable}{L{0.20\linewidth}Y}
\toprule
\textbf{层} & \textbf{典型检查对象}\\
\midrule
Retrieval & Recall@k、precision / relevance、MRR / ranking quality、permission-filter correctness\\
Generation & answer correctness、groundedness、citation correctness、refusal correctness、structured-output validity\\
Agent & task success、tool selection、tool argument accuracy、trajectory quality、unnecessary calls、loop / stop correctness、safety violations\\
System & P50 / P95 latency、token / tool cost、failure rate、recovery success、framework / model update regression\\
\bottomrule
\end{handbooktable}

任何 Model、Prompt、Retriever、Framework 或 MCP Server 更新，都应能跑 regression eval，而不是只靠人工试几次。

\subsection{Agent Eval 的基本单位：Task、Trial、Grader、Transcript}

Agent Eval 和单轮问答最大的区别是：结果会受到多步决策、工具、环境状态和随机性的共同影响。为了避免“跑一次看起来不错”这种假精确，至少把四个对象分开：

\begin{handbooktable}{L{0.20\linewidth}Y}
\toprule
\textbf{对象} & \textbf{含义}\\
\midrule
Task & 一个有明确输入、环境初始状态和 success criteria 的测试问题\\
Trial & 同一个 Task 的一次独立运行；因为 Agent path 有随机性，一个 Task 通常需要多个 Trials\\
Grader & 对 outcome 或 trajectory 做检查的代码、模型或人工规则；一个 Task 可以有多个 Graders\\
Transcript / Trace & 一次 Trial 的完整执行证据，包括模型 turns、tool calls、observations、state changes 与 final outcome\\
\bottomrule
\end{handbooktable}

因此一个有意义的结果通常是分布而不是单点：

$$
\text{Task Success Rate}
=\frac{\text{successful trials}}{\text{all comparable trials}}.
$$

如果只保存 final answer，就无法判断失败是路径偶然、Tool 环境异常，还是 Harness / Model 的稳定缺陷。

\subsection{Agentic Benchmark 里 Runtime 本身就是实验变量}

传统静态 benchmark 常把模型输出直接交给评分器；Agentic benchmark 则让模型在完整环境里安装依赖、运行测试、读写文件并迭代。因此 CPU、RAM、timeout、network、container policy、package cache 和并发策略都会改变实际可解问题集合。

Anthropic 2026 对 Terminal-Bench 2.0 的实验中，仅基础设施资源配置就造成最高约 6 个百分点的 score spread。这一数字只属于该实验，不能推广成固定阈值；但它证明了一个更一般的实验原则：

\begin{criterion}{Agent Eval 可比性}
比较两个模型或两个 Harness 前，必须同时固定并报告 Model Config、Prompt/Harness、Tool Versions、Environment Image、CPU/RAM、Timeout、Network Policy、Retry Policy 与 Trial Count。否则小幅 leaderboard 差异可能只是 Runtime 条件不同。
\end{criterion}

这也是为什么 Agent Eval 应同时保存 outcome、trajectory 与 environment metadata。

\subsection{LLMOps：Trace 只有进入评测与回放闭环才会产生改进}

“记录很多日志”不等于 LLMOps。一个可维护的 Agent 系统至少需要形成：

\begin{processchain}
  \processstage{Trace Runs}
  \processstage{Evaluate}
  \processstage{Diagnose}
  \processstage{Change Config / Code}
  \processstage{Replay / Regression}
  \processstage{Deploy / Roll Back}
\end{processchain}

其中：

\begin{itemize}
  \item \textbf{Trace} 保存足以重建行为的可观察事实：inputs/references、model/tool calls、tool results、state transitions、errors、latency、tokens、cost；不要求记录模型不可用的私有 chain-of-thought；
  \item \textbf{Evaluate} 使用 deterministic assertions、reference-based metrics、task-specific graders、必要时 LLM-as-judge 与 human review；
  \item \textbf{Diagnose} 必须定位 first failing layer，是 retrieval miss、bad routing、tool schema、permission、stop condition、model response 还是业务 side effect；
  \item \textbf{Change} 把 prompt、model、retrieval config、tool definition、memory policy 与 workflow topology 都作为可版本化配置，而不是线上随手改字符串；
  \item \textbf{Replay / Regression} 在固定 eval set 与历史失败 case 上重跑，证明新版本修复目标问题且没有制造新的回归；
  \item \textbf{Deploy} 对高风险系统保留 staged rollout、canary、rollback 与变更审计。
\end{itemize}

自动“自我改进”尤其需要这一层约束。Agent 可以提出新的 prompt、skill 或 memory rule，但不能因为一次运行看起来成功就直接把它提升为生产默认；先经过可重复 eval 与 review，才能把一次经验变成稳定系统行为。

Trace 本身也属于敏感数据资产。用户输入、Tool arguments/results、Memory、身份信息和 secrets 进入 tracing backend 前应执行最小化、redaction、retention 与 access-control policy。

\subsection{Cost 与 Latency 是多步系统的累积量}

一次 Agent 任务的总成本可以近似写成：

$$
C_{\text{total}}
=\sum_i C_{\text{model},i}
+\sum_j C_{\text{tool},j}
+C_{\text{retrieval}}
+C_{\text{infra}}.
$$

总延迟可以近似写成：

$$
T_{\text{total}}
\approx T_{\text{retrieval}}
+\sum_i T_{\text{model},i}
+\sum_j T_{\text{tool},j}
+T_{\text{network}}.
$$

更多 Context、更多 Agent iterations、更多 Tools 和更多 verification passes 都可能提高质量，也都会提高 cost / latency。是否值得只能通过实际任务评测决定。

\subsection{Failure Recovery：失败是运行时的一部分}

至少要设计：

\begin{itemize}
  \item LLM / Retriever timeout；
  \item MCP server unavailable；
  \item Tool validation / business error；
  \item partial write；
  \item duplicate retry；
  \item process crash；
  \item graph loop；
  \item stale checkpoint；
  \item malformed structured output。
\end{itemize}

读操作重试通常比较安全；写操作必须考虑 idempotency；破坏性写入通常需要 explicit confirmation 或 compensation。LangGraph checkpoint 可以恢复 workflow state，但不能自动替业务系统解决重复副作用。

\section{Complexity Ladder：什么时候不要继续增加 Agent 自主性}

\begin{handbooktable}{L{0.17\linewidth}L{0.27\linewidth}Y}
\toprule
\textbf{层级} & \textbf{结构} & \textbf{什么时候足够}\\
\midrule
Level 0 & 普通确定性代码 & 规则已知，输入输出关系可以直接写成程序\\
Level 1 & Single LLM Call & summarization、extraction、classification、one-shot generation / reasoning\\
Level 2 & Retrieval / RAG & 需要私有知识、大量文档、最新资料或 citation\\
Level 3 & Tool Calling & 模型需要读取或操作实时外部系统\\
Level 4 & Fixed Workflow / Graph & 有固定阶段、分支、retry、approval、checkpoint 和明确 state transition\\
Level 5 & Autonomous Agent & 下一步无法在设计时可靠确定，需要模型根据环境与 Tool observation 动态选动作\\
\bottomrule
\end{handbooktable}

MCP 不是 Level 6。它横切 Level 3--5，为外部 Context / Capability 提供标准化连接方式。Multi-Agent 也不是 Level 6；它是对某些可分解任务进行 context isolation、specialization 或 parallel scaling 的组织方式，可以出现在 Workflow 或 Autonomous Agent 内部。

\begin{criterion}{复杂度升级判据}
能用普通函数解决，就不要上 Agent；能用固定 Workflow 清楚表达控制流，就不要把每一步都交给模型自由决定。只有“下一步本身依赖开放环境和中间观察，设计时无法稳定穷举”时，更高自主性才有充分理由。
\end{criterion}

\section{一条可执行的实验路线}

\subsection{Lab A｜直接 Model API}

目标：理解最小 model runtime。

\begin{processchain}
  \processstage{OpenAI SDK}
  \processstage{Responses API}
  \processstage{Input}
  \processstage{Output}
  \processstage{Usage / Latency}
\end{processchain}

至少记录 model、input tokens、output tokens、latency 与 cost。

\subsection{Lab B｜Semantic Retrieval}

建立：Documents $\to$ Chunk $\to$ Embedding $\to$ Index $\to$ Search。至少比较 keyword vs semantic、small vs large chunk、with vs without metadata filter。准备一个小型 labeled query set，不用单个 query 宣布效果。

\subsection{Lab C｜RAG}

流程：Query $\to$ Retrieve Top-k $\to$ Build Context $\to$ Generate $\to$ Cite Sources。分别检查 retrieval hit、answer correctness、citation support，以及 evidence 不足时是否正确拒答。

\subsection{Lab D｜Function / Tool Calling}

实现一个 read-only lookup 与一个 state-changing operation，对比 schema、permission、confirmation、retry 和 audit log。这样才能真正看到 Tool Safety 的差异。

\subsection{Lab E｜LangChain Agent}

使用 \codepath{langchain.agents.create_agent} 组合 model + tools，重点记录 model decision $\to$ tool call $\to$ observation $\to$ next decision $\to$ stop 的完整 trajectory。

\subsection{Lab F｜LangGraph Workflow}

把同一任务改成有显式 State 的图：START $\to$ Retrieve $\to$ Analyze $\to$ Need More Evidence?；证据不足时回到 Retrieval / Tool，足够时进入 Generate $\to$ Verify $\to$ END。加入 checkpointer、iteration limit、interrupt / approval、deliberate failure 与 resume test。

\subsection{Lab G｜MCP Server}

使用当前 Python SDK v2：\code{MCPServer} $\to$ \code{@mcp.tool()} $\to$ \code{stdio}，再改成 remote Streamable HTTP。至少暴露一个 Tool、一个 Resource、一个 Prompt，避免把 MCP 误解成只有 Tool API。

\subsection{Lab H｜Production Evaluation}

最后主动注入 Normal Case、Edge Case、Prompt Injection、Retriever Miss、Tool Timeout、Tool Error、Duplicate Retry、Permission Denial、MCP Server Offline。系统只在 happy path 演示成功，还不能称为可靠。

\subsection{Lab I｜Long-term Memory Lifecycle}

建立一组跨 session 的小型记忆测试集，而不是只问一个“还记得我吗”。至少包含：稳定偏好、精确 ID、跨语言改写、时间更新、旧事实 supersede、应该忘记的敏感信息、两个用户之间的隔离，以及一条来自不可信网页/工具输出的恶意 instruction。分别记录：write decision、stored representation、retrieved top-k、final answer、latency、token/cost 与后续删除结果。再比较 always-in-context、FTS/BM25、vector/hybrid 与 temporal graph 中真正适合当前数据结构的方案。

\subsection{Lab J｜Skill、Tool 与 Harness Boundary}

选择一个重复但多步骤的任务，例如“初始化一个 RAG backend 并验证 health endpoint”。先只给 Agent 原始 shell/file Tools，再把稳定步骤写成 Skill。比较两次运行的 instruction repetition、错误率与 token usage。随后故意删除一个底层 Tool，验证 Skill 是否会因为“写了步骤”就 magically 获得执行能力；正确结果应该是 Skill 仍然只能指导，真正动作必须由 Tool/Runtime 提供。

\subsection{Lab K｜Loop、Graph 与 Hybrid}

用同一个业务目标实现三版：纯 Agent Loop、固定 Graph、Graph 中嵌入一个 Agent research/debug loop。注入一个未知外部问题和一个确定性审批步骤，观察哪部分真正需要模型临场决策。记录 stop correctness、unnecessary tool calls、side-effect safety、latency 与 recovery complexity，避免只比较最终答案。

\subsection{Lab L｜Orchestrator--Worker Deep Research}

准备一个可自然拆成 3--5 个独立方向的研究任务。Single-Agent 作为 baseline；Multi-Agent 版本由 Orchestrator 明确 worker objective、tools、budget、output schema，Workers 把 findings + sources 写入 artifacts，再由 Synthesizer 汇总。比较 coverage、duplicate work、handoff loss、wall-clock latency、total tokens/cost 与 citation correctness。如果任务本来高度串行，应看到多 Agent 的协调成本可能大于收益。

\subsection{Lab M｜Session、Context 与 Crash Recovery}

构造一个会产生多个 artifacts 的长任务。把完整 event log / task state 存在 Context Window 外部；运行若干步骤后主动杀死 Harness 进程或丢弃当前 Context，再从 durable Session 恢复。检查：是否能定位最后一个已提交副作用、是否会重复写入、是否能重新选择必要历史而不是把全部 log 再塞回 Context。随后对比“只有 conversation summary、没有原始 event log”的版本，观察哪些错误变得不可审计或不可恢复。

\subsection{Lab N｜Agent Eval 的重复试验与 Runtime 控制}

为同一 Task 至少运行多次 independent trials，并同时记录 outcome grader、trajectory grader、token/cost 和 environment metadata。再只改变一个 Runtime 变量，例如 timeout 或 RAM limit，观察 success distribution 是否变化。这个实验的目标不是追求某个 leaderboard 数字，而是亲自验证：Agentic Eval 测量的是 \textbf{Model + Harness + Tools + Environment} 的组合行为。

\section{最终压缩：十四个问题复原整套系统}

遇到任何 LLM / Agent 系统，依次问：

\begin{enumerate}
  \item \textbf{模型当前真正看见什么？}——Context；
  \item \textbf{当前证据从哪里来？}——Retrieval / Files / DB / API；
  \item \textbf{证据怎样进入模型？}——Context Construction / RAG；
  \item \textbf{哪些信息要跨会话保留，何时写入、更新和取回？}——Long-term Memory；
  \item \textbf{模型只能回答，还是可以行动？}——Tool Calling；
  \item \textbf{重复任务的“怎样做”是否需要独立沉淀？}——Skill / Procedural Knowledge；
  \item \textbf{谁把 Context、Tools、Loop、State、Memory、Permissions 与 Trace 组装成可运行代理？}——Agent Harness；
  \item \textbf{Session 与这一轮 Context 是否已经分开？}——Durable Event Log / Context View；
  \item \textbf{下一步是开放探索还是已知流程？}——Loop / Graph / Hybrid；
  \item \textbf{是否存在真正独立、值得隔离或并行的子任务？}——Subagent / Multi-Agent；
  \item \textbf{当前运行状态在哪里保存，失败后怎样恢复？}——State / Checkpoint / Artifacts；
  \item \textbf{真正执行副作用的环境与模型是否隔离？}——Sandbox / Permission / Credential / Egress Boundary；
  \item \textbf{外部能力怎样标准化接入？}——MCP / Tool Contract；
  \item \textbf{怎样证明修改后的系统真的更可靠？}——Repeated Trials / Trace / Evaluation / Security / LLMOps Regression。
\end{enumerate}

\begin{mentalmodel}{一页压缩}
Context 决定模型这一轮能看见什么，而 Durable Session 保存 Context Window 外部可恢复的任务历史；Retrieval / RAG 决定外部证据怎样进入；Long-term Memory 负责跨任务信息的持久化与再调用；Skills 沉淀重复流程，Tools 提供实际动作；Agent Harness 把这些能力连同状态、权限、循环和观测组织成可运行系统，并把 Model/Decision 与 Sandbox/Side Effects 解耦；Loop 处理下一步未知的开放任务，Graph 固定已知业务骨架，Multi-Agent 只在任务可分解或需要独立 evaluator 时扩展执行上下文；MCP 标准化外部能力接口；Containment 限制 Blast Radius，Repeated-trial Evaluation 与 LLMOps Regression 决定系统能否长期可靠演进。
\end{mentalmodel}

\section{当前技术快照与官方资料}

本节只锁定快速变化的 API / protocol 行为，机制主干不依赖具体厂商版本。当前快照日期为 2026-08-22。

\begin{itemize}
  \item OpenAI Models / GPT-5.6：\url{https://developers.openai.com/api/docs/models}
  \item OpenAI GPT-5.6 Sol：\url{https://developers.openai.com/api/docs/models/gpt-5.6-sol}
  \item OpenAI Responses API：\url{https://developers.openai.com/api/reference/resources/responses}
  \item OpenAI Vector Stores：\url{https://platform.openai.com/docs/api-reference/vector-stores}
  \item OpenAI Agents SDK：\url{https://openai.github.io/openai-agents-python/}
  \item OpenAI Agents SDK Tracing / Guardrails：\url{https://openai.github.io/openai-agents-python/tracing/}，\url{https://openai.github.io/openai-agents-python/guardrails/}
  \item LangChain Agents：\url{https://docs.langchain.com/oss/python/langchain/agents}
  \item LangGraph Workflows / Agents：\url{https://docs.langchain.com/oss/python/langgraph/workflows-agents}
  \item LangGraph Persistence：\url{https://docs.langchain.com/oss/python/langgraph/persistence}
  \item LangChain Deep Agents / Agent Harness：\url{https://docs.langchain.com/oss/python/deepagents/overview}
  \item Anthropic Agent Skills：\url{https://www.anthropic.com/engineering/equipping-agents-for-the-real-world-with-agent-skills}
  \item Anthropic Multi-Agent Research：\url{https://www.anthropic.com/engineering/multi-agent-research-system}
  \item Anthropic Long-running Agent Harness：\url{https://www.anthropic.com/engineering/effective-harnesses-for-long-running-agents}
  \item Anthropic Harness Design for Long-running Apps：\url{https://www.anthropic.com/engineering/harness-design-long-running-apps}
  \item Anthropic Managed Agents / Brain--Hands--Session：\url{https://www.anthropic.com/engineering/managed-agents}
  \item Anthropic Agent Containment：\url{https://www.anthropic.com/engineering/how-we-contain-claude}
  \item Anthropic Agent Evals / Infrastructure Noise：\url{https://www.anthropic.com/engineering/demystifying-evals-for-ai-agents}，\url{https://www.anthropic.com/engineering/infrastructure-noise}
  \item MCP 2026-07-28：\url{https://blog.modelcontextprotocol.io/posts/2026-07-28/}
  \item MCP Python SDK v2：\url{https://github.com/modelcontextprotocol/python-sdk}
  \item pgvector：\url{https://github.com/pgvector/pgvector}
  \item Supabase Vector Columns / pgvector：\url{https://supabase.com/docs/guides/ai/vector-columns}
  \item Hermes Agent Persistent Memory：\url{https://github.com/NousResearch/hermes-agent}
  \item Waku Agent Memory Architecture：\url{https://github.com/ShenSeanChen/waku-agent}
  \item LangMem：\url{https://langchain-ai.github.io/langmem/}
  \item Mem0 OSS v3 Migration / Graph Memory：\url{https://github.com/mem0ai/mem0}
  \item Graphiti / Zep Temporal Knowledge Graph：\url{https://github.com/getzep/graphiti}
\end{itemize}

\end{document}
